\documentclass[11pt,a4paper]{article}

\usepackage[utf8]{inputenc}
\usepackage[T1]{fontenc}
\usepackage[margin=1in]{geometry}
\usepackage{amsmath, amssymb, amsthm, mathtools}
\usepackage{graphicx}
\usepackage{caption}
\usepackage{subcaption}
\usepackage{algorithm}
\usepackage{algpseudocode}
\usepackage{enumitem}
\usepackage{xcolor}
\usepackage{hyperref}
\usepackage{cleveref}
\usepackage{booktabs}

% --- Theorem environments (each with its own counter so cleveref works) ---
\theoremstyle{plain}
\newtheorem{theorem}{Theorem}
\newtheorem{proposition}{Proposition}
\newtheorem{lemma}{Lemma}
\newtheorem{corollary}{Corollary}

\theoremstyle{definition}
\newtheorem{definition}{Definition}
\newtheorem{assumption}{Assumption}
\newtheorem{remark}{Remark}

\crefname{theorem}{Theorem}{Theorems}
\Crefname{theorem}{Theorem}{Theorems}
\crefname{proposition}{Proposition}{Propositions}
\Crefname{proposition}{Proposition}{Propositions}
\crefname{lemma}{Lemma}{Lemmas}
\Crefname{lemma}{Lemma}{Lemmas}
\crefname{corollary}{Corollary}{Corollaries}
\Crefname{corollary}{Corollary}{Corollaries}
\crefname{definition}{Definition}{Definitions}
\Crefname{definition}{Definition}{Definitions}
\crefname{assumption}{Assumption}{Assumptions}
\Crefname{assumption}{Assumption}{Assumptions}
\crefname{remark}{Remark}{Remarks}
\Crefname{remark}{Remark}{Remarks}

%-----------------------------------------------------------------
% Macros
%-----------------------------------------------------------------
\newcommand{\R}{\mathbb{R}}
\newcommand{\N}{\mathbb{N}}
\newcommand{\E}{\mathbb{E}}
\newcommand{\norm}[1]{\left\|#1\right\|}
\newcommand{\inner}[2]{\left\langle #1,#2 \right\rangle}
\newcommand{\eps}{\varepsilon}
\newcommand{\epsg}{\eps_{g}}
\newcommand{\epsH}{\eps_{H}}
\newcommand{\grad}{\nabla}
\newcommand{\hess}{\nabla^{2}}
\newcommand{\xstar}{x^{\star}}
\newcommand{\fstar}{f^{\star}}
\newcommand{\flow}{f_{\mathrm{low}}}
\newcommand{\Lone}{L_{1}}
\newcommand{\Ltwo}{L_{2}}
\newcommand{\setN}{\mathcal{N}}
\newcommand{\setR}{\mathcal{R}}
\newcommand{\cN}{c_{\setN}}
\newcommand{\cR}{c_{\setR}}
\newcommand{\Tp}{T_{\mathrm{perturb}}}
\newcommand{\Fthresh}{F_{\mathrm{thresh}}}
\newcommand{\tnoise}{t_{\mathrm{noise}}}
\newcommand{\snap}{\tilde{x}}
\newcommand{\fsnap}{\hat{f}}
\renewcommand{\d}{d}
\newcommand{\vmin}{v_{\min}}
\newcommand{\lmin}{\lambda_{\min}}
\newcommand{\sdel}{s_{\delta}}
\newcommand{\glow}{\underline{g}}
\renewcommand{\Pr}{\mathbb{P}}
\newcommand{\indic}[1]{\mathbf{1}_{\{#1\}}}
\newcommand{\proofpara}[1]{\par\medskip\noindent\textbf{#1}\quad}
\newcommand{\defeq}{\mathrel{\vcentcolon=}}
% --- Macros for the complexity analysis (Step IX) ---
\newcommand{\LL}{\mathcal{L}}                
\newcommand{\Rstar}{R^{*}}
\newcommand{\gup}{\bar g}
\newcommand{\bbeta}{\bar\beta}
\newcommand{\Niter}{N_{\mathrm{iter}}}
\newcommand{\Ntrig}{N_{\mathrm{trig}}}
\newcommand{\serr}{\sigma_{\mathrm{err}}}
\newcommand{\smom}{\sigma_{\mathrm{mom}}}
\newcommand{\tstar}{t_\star}
\newcommand{\thetamin}{\theta''_{\min}}
\newcommand{\amin}{\alpha_{\min}}

\title{Non-linear Conjugate Gradient With Restart Also Escapes Saddle Points\\
\large Working document --- full proofs and rationale}
\author{Benjamin S\'eguin, Cl\'ement Royer \\ \small LAMSADE}
\date{\today}

\begin{document}
\maketitle

\begin{abstract}
We study a perturbed, restarted variant of the non-linear conjugate gradient (NCG)
method and show that, like perturbed gradient descent, it escapes strict saddle
points and converges to an approximate second-order stationary point. The analysis
combines per-step decrease guarantees, an improve-or-localize argument valid for both
restarted and non-restarted iterations, a linearization of the dynamics around the
snapshot point, an anti-concentration estimate for the random perturbation, an
induction maintaining three coupled invariants over the escape window, and a
contradiction argument converting geometric growth into guaranteed function decrease.
A final section resolves the resulting constraint system exactly and derives the
complexity bound, together with its current limitations.
\end{abstract}

\tableofcontents

% ============================================================
\section{Introduction}\label{sec:intro}
We follow the framework of \cite{pmlr-v70-jin17a} and \cite{pmlr-v75-jin18a} and
adapt it to the non-linear conjugate gradient method with restart, in order to show
that this method also escapes strict saddle points and produces an approximate
second-order stationary point. The starting point for the descent estimates is the
NCG complexity analysis of \cite{RoyerChan2022}; the saddle-avoidance
mechanism is borrowed, at a high level, from the perturbed gradient descent analyses of
\cite{pmlr-v70-jin17a,pmlr-v75-jin18a} and from the dynamical-systems viewpoint of
\cite{lee2019}. Our setup is comparable to the one of~\cite{pmlr-v75-jin18a}, since
NCG is a momentum method, just like accelerated gradient descent (AGD).
However, we depart from their proof technique by studying a \emph{single trajectory},
whereas they study two coupled trajectories of the same process and conclude from
the growth of their difference.

The document is organized around the proof story of \Cref{sec:roadmap}: nine steps,
each with a one-paragraph rationale (\emph{what} it establishes, \emph{why} the
argument needs it, and \emph{where} it is consumed), followed by the corresponding
section with full proofs.

\section{Preliminaries}\label{sec:prelim}
In this section, we introduce our notation, and present the definitions, assumptions,
and elementary consequences used throughout this document.

\subsection{Notation}
For vectors, we use
$\norm{\cdot}$ to denote the $\ell_2$-norm. We use $\lambda_{\max}(\cdot)$ and
$\lambda_{\min}(\cdot)$ to denote the largest and smallest eigenvalues respectively.
For a function $f \colon \R^{d} \to \R$, we use $\grad f(\cdot)$ and $\hess f(\cdot)$ to
denote its gradient and Hessian, and $\fstar$ to denote the global minimum of $f$. We
let $\mathcal{B}^{d}(x, r)$ denote the $d$-dimensional ball centered at $x$ with radius
$r$, and $\mathcal{U}(S)$ the uniform distribution over a set $S$.

\subsection{Setting}
We consider the unconstrained minimization problem
\begin{equation}\label{eq:problem}
\min_{x \in \R^{d}} f(x),
\end{equation}
where $f \colon \R^{d} \to \R$ is nonconvex and whose saddle points are strict.

\begin{definition}[Strict saddle point]\label{def:strictsaddle}
A strict saddle point is a saddle point $\xstar$ that has at least one direction of
negative curvature, that is,
\begin{equation}
\lambda_{\min}\!\big(\hess f(\xstar)\big) < 0 .
\end{equation}
\end{definition}

\noindent The algorithm certifies approximate second-order stationarity in the following sense.

\begin{definition}[Approximate second-order stationary point]\label{def:2osp}
Given tolerances $\epsg, \epsH > 0$, a point $x \in \R^{d}$ is an
$(\epsg, \epsH)$-second-order stationary point (SOSP) if
\begin{equation}\label{eq:2osp}
\norm{\grad f(x)} \leq \epsg
\qquad\text{and}\qquad
\lambda_{\min}\!\big(\hess f(x)\big) \geq -\epsH .
\end{equation}
\end{definition}

\noindent We make the following standing assumptions.

\begin{assumption}[Gradient Lipschitz]\label{ass:gradlip}
The function $f$ is continuously differentiable on $\R^{d}$, and its gradient $\grad f$
is $\Lone$-Lipschitz continuous, with $\Lone > 0$.
\end{assumption}

\begin{assumption}[Hessian Lipschitz]\label{ass:hesslip}
The function $f$ is twice continuously differentiable on $\R^{d}$, and its Hessian
$\hess f$ is $\Ltwo$-Lipschitz continuous, with $\Ltwo > 0$.
\end{assumption}

\subsection{Elementary consequences of the assumptions}

Both consequences below are used repeatedly; we record them once with full proofs.

\begin{lemma}[Descent lemma]\label{lem:descent}
Under \Cref{ass:gradlip}, for all $x, y \in \R^d$,
\begin{equation}\label{eq:descent}
f(y) \;\le\; f(x) + \grad f(x)^{\top}(y-x) + \frac{\Lone}{2}\norm{y-x}^{2}.
\end{equation}
\end{lemma}

\begin{proof}
Let $\varphi(t) := f\bigl(x + t(y-x)\bigr)$ for $t\in[0,1]$, so that
$\varphi'(t) = \grad f\bigl(x+t(y-x)\bigr)^{\top}(y-x)$. By the fundamental theorem of
calculus,
\[
f(y) - f(x)
= \int_{0}^{1}\varphi'(t)\,dt
= \grad f(x)^{\top}(y-x)
  + \int_{0}^{1}\bigl(\grad f(x+t(y-x)) - \grad f(x)\bigr)^{\top}(y-x)\,dt .
\]
By Cauchy--Schwarz and \Cref{ass:gradlip}, the integrand is bounded by
$\Lone t \norm{y-x}^{2}$, and $\int_{0}^{1}\Lone t\,dt = \Lone/2$, giving
\eqref{eq:descent}.
\end{proof}

\begin{lemma}[Spectral bound on the Hessian]\label{lem:hessbound}
Under \Cref{ass:gradlip,ass:hesslip}, for every $x\in\R^{d}$,
$\norm{\hess f(x)} \le \Lone$; in particular
$\lambda_{\min}(\hess f(x)) \ge -\Lone$ everywhere.
\end{lemma}

\begin{proof}
Fix $x$ and a unit vector $v$. By \Cref{ass:gradlip},
$\norm{\grad f(x+tv) - \grad f(x)} \le \Lone |t|$ for all $t$, so
\[
\norm{\hess f(x)\,v}
= \lim_{t\to 0}\frac{\norm{\grad f(x+tv)-\grad f(x)}}{|t|}
\le \Lone .
\]
Taking the supremum over unit $v$ gives $\norm{\hess f(x)}\le\Lone$; since
$\hess f(x)$ is symmetric, all its eigenvalues lie in $[-\Lone,\Lone]$.
\end{proof}

\begin{remark}[Meaningful range of $\epsH$]\label{rem:epsH-range}
By \Cref{lem:hessbound}, the curvature condition
$\lambda_{\min}(\hess f(x)) \le -\epsH$ can hold at \emph{some} point only if
$\epsH \le \Lone$. For $\epsH > \Lone$ the perturbation trigger of
\Cref{algo:PNCG} can never fire and the SOSP certificate \eqref{eq:2osp} holds
automatically at every $\epsg$-stationary point. This elementary observation
becomes important when solving the constraints system.
% reading the regime restriction of Step~IX; see
% \Cref{rem:vacuity}.
\end{remark}

% ============================================================
\section{The Perturbed Restarted NCG Algorithm}\label{sec:algo}

We present a perturbed variant of the NCG method with
restart, analogous to the perturbed gradient descent of \cite{pmlr-v70-jin17a}. The
algorithm augments the usual NCG iteration with four ingredients absent from the plain
method: a perturbation radius $r$, an escape window $\Tp$, a function-value progress
threshold $\Fthresh$, and a target failure probability $\delta$. The idea is
straightforward: when the norm of the current gradient is small ($\norm{g_k} \leq
\epsg$), the current point is a candidate SOSP, and a random
perturbation is injected to test whether it is a SOSP or a
strict saddle point.

\begin{algorithm}[H]
\caption{Perturbed Restarted Non-linear Conjugate Gradient (PR-NCG)}\label{algo:PNCG}
\begin{algorithmic}[1]
\Require $x_{0} \in \R^{d}$; NCG parameters $\sigma \in (0,1]$, $\kappa \geq 1$,
  $p \geq 0$, $q \geq 0$; tolerances $\epsg > 0$, $\epsH > 0$; perturbation radius
  $r > 0$, escape window $\Tp \in \N$, progress threshold $\Fthresh > 0$, failure
  probability $\delta \in (0,1)$.
\State $g_{0} \gets \grad f(x_{0})$,\quad $d_{0} \gets -g_{0}$
\State $\tnoise \gets -\Tp - 1$ \Comment{allows an immediate perturbation}
\State $\snap \gets x_{0}$,\quad $\fsnap \gets f(x_{0})$ \Comment{last snapshot}
\For{$k = 0, 1, 2, \dots$}
  \State $g_{k} \gets \grad f(x_{k})$
  \If{$\norm{g_{k}} \leq \epsg$ \textbf{ and } $k - \tnoise > \Tp$}
       \Comment{perturbation trigger (replaces NCD)}
    \State $\snap \gets x_{k}$,\quad $\fsnap \gets f(x_{k})$ \Comment{save snapshot}
    \State sample $\xi \sim \mathcal{U}\!\big(\mathcal{B}^{d}(0, r)\big)$
    \State $x_{k} \gets x_{k} + \xi$
    \State $g_{k} \gets \grad f(x_{k})$ \Comment{recompute after perturbation}
    \State $d_{k} \gets -g_{k}$ \Comment{reset momentum}
    \State $\tnoise \gets k$
  \EndIf
  \If{$k - \tnoise = \Tp$ \textbf{ and } $f(x_{k}) - \fsnap > -\Fthresh$}
       \Comment{termination check}
    \State \Return $\snap$ as an approximate $(\epsg, \epsH)$-2OSP
  \EndIf
  \Statex \quad\textit{/* NCG step */}
    \If{$d_{k} = -g_{k}$} \Comment{restarted direction}
    \State $\alpha_{k} \gets \alpha_{\setR}^{\ast} = 1/\Lone$
  \Else
    \State $\alpha_{k} \gets \alpha^{\ast}_{\setN}
    = \dfrac{\sigma}{\Lone \kappa^{2}}\,\norm{g_{k}}^{\,1+p-2q}$
  \EndIf
  \State $x_{k+1} \gets x_{k} + \alpha_{k} d_{k}$
  \State $g_{k+1} \gets \grad f(x_{k+1})$
  \State choose $\beta_{k+1}$
  \State $d_{k+1} \gets -g_{k+1} + \beta_{k+1} d_{k}$
  \If{$g_{k+1}^{\top} d_{k+1} \geq -\sigma \norm{g_{k+1}}^{\,1+p}$
       \textbf{ or } $\norm{d_{k+1}} \geq \kappa \norm{g_{k+1}}^{\,q}$}
       \Comment{restart condition}
    \State $d_{k+1} \gets -g_{k+1}$
  \EndIf
\EndFor
\end{algorithmic}
\end{algorithm}

We write $\setN$ for the set of non-restarted iterations (the kept NCG direction
satisfies $g_{k}^{\top}d_{k} < -\sigma\norm{g_k}^{1+p}$ and
$\norm{d_k} < \kappa\norm{g_k}^{q}$, by the complement of the restart test) and
$\setR$ for the restarted ones ($d_k = -g_k$).

\begin{assumption}[Exponent matching]\label{ass:pq}
Throughout the escape analysis we take $1+p = 2q$ with $q = 1$, i.e.\
$p = q = 1$. Under this assumption the non-restarted stepsize is the
\emph{constant}
\[
\alpha^{\ast}_{\setN} \;=\; \frac{\sigma}{\Lone\kappa^{2}},
\qquad\text{with}\qquad
\alpha^{\ast}_{\setN}\,\Lone \;=\; \frac{\sigma}{\kappa^{2}} \;\le\; 1 ,
\]
since $\sigma\le 1\le\kappa^{2}$. We write
\begin{equation}\label{eq:amin}
\amin \;:=\; \frac{\sigma}{\Lone\kappa^{2}},
\qquad\text{so that every stepsize satisfies}\qquad
\alpha_{k}\in\Bigl\{\amin,\ \tfrac{1}{\Lone}\Bigr\}\subset
\Bigl[\amin,\ \tfrac{1}{\Lone}\Bigr],
\quad \alpha_{k}\Lone\le 1 .
\end{equation}
The restart test then reads $\norm{d_{k+1}}\ge\kappa\norm{g_{k+1}}$ and
$g_{k+1}^{\top}d_{k+1}\ge-\sigma\norm{g_{k+1}}^{2}$.
\end{assumption}

\begin{remark}[Changelog: removal of the in-window stepsize safeguard]
\label{rem:safeguard-changelog}
% NOTE (flagged correction, kept for the record):
An earlier version of \Cref{algo:PNCG} had the stepsize at $\alpha=1/\Lone$
inside the escape window, for \emph{all} directions. This created a third
iteration type: NCG direction with $\alpha=1/\Lone$, which was covered by neither
case of \Cref{lem:improve}, and for which the per-step decrease can
fail: the descent lemma only gives
$f(x_{k})-f(x_{k+1})\ge\tfrac{\sigma}{\Lone}\norm{g_{k}}^{1+p}
-\tfrac{\kappa^{2}}{2\Lone}\norm{g_{k}}^{2q}$, whose sign is uncontrolled in
general. Since \Cref{lem:localization} applies \Cref{lem:improve} to in-window
steps, this was a gap.

\emph{Resolution adopted.} Under \Cref{ass:pq} the stepsize
$\alpha^{\ast}_{\setN}$ is constant, so the safeguard is unnecessary: the
window runs the algorithm's ordinary stepsize rule, every in-window step is
one of the two cases of \Cref{lem:improve} (restarted with $1/\Lone$, or
non-restarted with $\alpha^{\ast}_{\setN}$), and the uniform constant
$2/\Lone$ applies verbatim. The escape recurrence becomes variable-stepsize
with $\alpha_{k}\in[\amin,1/\Lone]$; the growth rate is controlled by
$\amin$, costing the absolute factor $\kappa^{2}/\sigma$ in $T$ and
$\kappa^{4}/\sigma^{2}$ in $\Niter$, with no change to the
$(\epsH,\d,\delta)$ dependence.

\emph{Alternatives not adopted.} (a)~Keeping $\alpha=1/\Lone$ in-window and
requiring $\kappa^{2}<2\sigma$ also restores a per-step bound, with the
non-uniform constant $2c_{w}/\Lone$, $c_{w}=\kappa^{2}/(2\sigma-\kappa^{2})$;
valid (localization only needs the maximum of the per-step constants), but it
confines $(\sigma,\kappa)$ to $\kappa\in[1,\sqrt{2\sigma})$.
(b)~For general $(p,q)$, the safeguard
$\alpha_{k}=\min\{1/\Lone,\alpha^{\ast}_{\setN}\}$ restores the bound but
makes the in-window growth rate depend on $\glow^{\,1+p-2q}$, degrading the
final rate; not pursued.
% TODO: if the general-(p,q) extension is ever needed, the natural route is a
% Hamiltonian/potential argument a la Jin-Netrapalli-Jordan rather than (b).
\end{remark}

\subsection{Perturbation as a saddle detector}\label{sub:detector}
The perturbation plays the role of a strict-saddle detector. The intuition is that when
the gradient is small, we are at a candidate SOSP, and we perturb to get one of the
following two conclusions:
\begin{itemize}[leftmargin=*]
  \item If we perturb a point close to a strict saddle and let the algorithm run, the
  iterate moves a lot, because the negative-curvature direction amplifies the
  perturbation exponentially, and the function value drops significantly.
  \item If we perturb a point close to a local minimizer (a genuine SOSP) and let the
  algorithm run, the iterate stays roughly the same and the function value stagnates.
\end{itemize}
Hence the test ``did $f$ drop substantially after the perturbation?'' answers the
question ``was I at a strict saddle?''. This permits us to certify convergence to an SOSP
without any explicit computation of the Hessian matrix:
condition $\norm{g_{k}} \leq \epsg$ is trivial to check, while the curvature condition
$\lambda_{\min}(\hess f(x_{k})) \geq -\epsH$ is certified \emph{without} an explicit
Hessian computation, through the perturbation test described above.

% ============================================================
\section{Proof Story (How to Read This Document)}\label{sec:roadmap}

The argument decomposes into nine steps. Steps~I--II are generic descent facts about
the algorithm; Steps~III--VI build the toolbox for the near-saddle analysis; Step~VII
is the heart of the proof (the induction); Step~VIII converts it into the escape
theorem; Step~IX resolves the constants and counts iterations. For each step we
record \emph{what} it establishes, \emph{why} the argument needs it, and
\emph{where} it is proved.

\paragraph{Step I --- Decrease lemmas (\Cref{sec:decrease}).}
Basic: each regime (restarted, non-restarted) gives a per-step decrease guarantee,
with the stepsize $\alpha$ chosen to make the decrease optimal in that regime.
\emph{Why:} these are the raw material of every function-value accounting in the
document, in particular of Step~II.

\paragraph{Step II --- Improve or localize (\Cref{sec:improve}).}
We prove
\[
\sum_{k=0}^{T-1}\norm{x_{k+1}-x_{k}}^{2}
\;\le\; \frac{2}{\Lone}\bigl(f(x_{0})-f(x_{T})\bigr).
\]
There are two readings of this inequality.
\emph{(1) Improve:} if the iterates have moved a lot (left-hand side large), the
function value must have dropped a lot (right-hand side large).
\emph{(2) Localize:} if the function value has not dropped much (small right-hand
side), the iterates cannot have moved much (small left-hand side).
This is a dichotomy: at any moment, either we have improved by a meaningful amount,
or we are stuck near where we started.
\emph{Why:} the entire escape proof works inside the \emph{localize} branch and
derives a contradiction there; the dichotomy then forces the \emph{improve} branch.

\paragraph{Step III --- Linearization around $\snap$ (\Cref{sub:linearization}).}
Here we linearize in order to perform a spectral analysis (i.e.\ get a linear system
and compute eigenvalues). The whole idea is to show that the system has an eigenvalue
$\mu_{+}>1$: since the element of study is $u_{k}=x_{k}-\snap$, if $\mu_{+}>1$, the
iterate escapes the stationary point along the negative-curvature direction. We let
$c_{k}:=\inner{u_{k}}{\vmin}$ and study its recurrence to get geometric growth. The
recurrence is
\[
|c_{k+1}| \;\ge\; (1+\alpha\epsH)\,|c_{k}|
\;-\; \alpha\,\beta_{k}|\eta_{k}| \;-\; \alpha\,\epsg \;-\; \alpha\norm{\tilde e_{k}},
\]
with $\eta_{k}=\inner{d_{k-1}}{\vmin}$ a momentum term and
$\tilde e_{k}=\grad f(x_{k})-g_{\snap}-H(x_{k}-\snap)$ an error term. The rest of the
proof exists to control the other terms in the recurrence (momentum and quadratic
error) and to show that escaping along negative curvature has high probability.

\paragraph{Step IV --- Quadratic model error (\Cref{sub:graderror-section}).}
Straightforward: bound $\norm{\tilde e_{k}}$ above, which corresponds to the error
between the real gradient and the quadratically approximated one.
\emph{Why:} the recurrence of Step~III is only useful if this error cannot cancel
the geometric push.

\paragraph{Step V --- Probabilistic anti-concentration (\Cref{sub:anticonc}).}
Also a straightforward step. The idea is to show that the perturbation has a
non-negligible component along the negative-curvature direction, with high
probability. \emph{Why:} the geometric growth $|c_{k}|\gtrsim\mu_{+}^{k}|c_{0}|$ is
vacuous if $|c_{0}|$ is zero or tiny; this step supplies the initial signal
$|c_{0}|\ge\sdel$.

\paragraph{Step VI --- Lower bound on $\glow$ (Fletcher--Reeves specific)
(\Cref{sub:gfloor}).}
This step exists to control the momentum term. Since
$\beta^{\mathrm{FR}}_{k}=\norm{g_{k}}^{2}/\norm{g_{k-1}}^{2}$, we have to bound
$\beta^{\mathrm{FR}}\le \gup^{2}/\glow^{2}$, which requires a \emph{lower} bound on
gradient norms. Any other $\beta$ formula would mean redoing this step; heavy ball or
Nesterov do not need it, because their momentum coefficient is constant.
\emph{Why:} this is the FR-specific price of the analysis, and the source of the
final regime restriction.

\paragraph{Step VII --- Induction for geometric growth (\Cref{ssec:induction}).}
We want $|c_{k}|\ge(1+\gamma\alpha\epsH)^{k}|c_{0}|$. The issue is that the obstacle
terms in the recurrence depend on the past:
$\norm{\tilde e_{k}}\le\tfrac{\Ltwo}{2}\norm{u_{k}}^{2}$ depends on
$u_{k}=x_{k}-\snap$; $\beta_{k}$ depends on $\norm{g_{k-1}}$;
$|\eta_{k}|=|\inner{d_{k-1}}{\vmin}|$ depends on previous directions.
The solution is an induction, which lets us maintain a set of properties
simultaneously up to time $k$ and prove them at time $k+1$. From the recurrence, the
three things we need are: (1) how far the iterate is
($\norm{x_{k}-\snap}\le\dots$); (2) how big the gradient is
($\norm{g_{k}}\ge\glow$); (3) how big $|c_{k}|$ is ($|c_{k}|\ge\dots|c_{0}|$).
The whole idea is then driven by improve-or-localize:
\emph{(a)} the function value drops by at least $F$ somewhere in $\{0,\dots,T\}$:
this is what we try to prove;
\emph{(b)} the function value never drops by that much, in which case Step~II tells
us the iterates stay close to $x_{0}$
($\norm{x_{k}-x_{0}}\le S$, $S=\sqrt{2FT/\Lone}$).
We work in case~(b) and derive a contradiction, so we must be in case~(a). Case~(b)
is the bad case and we formalize it by the localization event
$\mathcal{L}_{T}=\{\max_{0\le k\le T}(f(x_{0})-f(x_{k}))<F\}$.

\paragraph{Step VIII --- Escape theorem (\Cref{sub:escape}).}
Using invariant~(3) of the induction, show
$|c_{T}|\ge(1+\gamma\alpha\epsH)^{T}\sdel$, and using invariant~(1),
$|c_{T}|\le\norm{u_{T}}\le S+\Rstar$. For a suitable $T$,
$(1+\gamma\alpha\epsH)^{T}\sdel > S+\Rstar$: we have the contradiction
$|c_{T}|\le S+\Rstar$ and $|c_{T}|>S+\Rstar$. Hence
$\mathcal{G}_{0}\cap\mathcal{L}_{T}=\emptyset$ (since we have been conditioning on
$\mathcal{G}_{0}$). This means we are in the first case of improve-or-localize, and
the function value drops by at least $F$ within the escape window with high
probability.

% \paragraph{Step IX --- Complexity (\Cref{sec:setting-cs}--\Cref{sec:main}).}
% The escape theorem holds \emph{provided} the parameters
% $(S,T,F,\theta,\theta',\theta'',\gamma)$ satisfy a system of constraints accumulated
% along Steps~III--VIII. Step~IX resolves this system exactly: it derives the two
% envelope bounds on $S$, optimizes the budget splits, computes
% $S$, $T$, $F$, and the iteration count, and closes the union bound over windows.
% The current limitations of the result (regime restriction and its strength) are
% discussed in \Cref{rem:vacuity}.

% ============================================================
\section{Step I --- Decrease Lemmas}\label{sec:decrease}

\paragraph{Rationale.}
Each iteration type comes with a stepsize chosen to make the guaranteed decrease
optimal for that type. The two lemmas below are the per-step guarantees; they feed
the improve-or-localize inequality of Step~II, and (through it) the entire
function-value accounting of the escape proof.

\begin{lemma}[Decrease on non-restarted iterations]\label{lem:decreaseN}
For every non-restarted iteration $k \in \setN$ (so that
$g_{k}^{\top}d_{k} \le -\sigma\norm{g_{k}}^{1+p}$,
$\norm{d_{k}} \le \kappa\norm{g_{k}}^{q}$, and
$\alpha_{k} = \alpha^{\ast}_{\setN} = \tfrac{\sigma}{\Lone\kappa^{2}}
\norm{g_{k}}^{1+p-2q}$),
\begin{equation}\label{eq:decreaseN}
f(x_{k}) - f(x_{k+1}) \;\geq\; \cN \, \norm{g_{k}}^{\,2(1+p-q)},
\qquad
\cN = \frac{\sigma^{2}}{2\Lone \kappa^{2}} .
\end{equation}
\end{lemma}

\begin{proof}
This is the fixed-stepsize NCG decrease guarantee of \cite{RoyerChan2022}; we
reproduce the computation for completeness. Apply the descent lemma
(\Cref{lem:descent}) at $x = x_{k}$, $y = x_{k+1} = x_{k}+\alpha_{k}d_{k}$:
\[
f(x_{k+1})
\;\le\; f(x_{k}) + \alpha_{k}\, g_{k}^{\top}d_{k}
        + \frac{\Lone}{2}\,\alpha_{k}^{2}\norm{d_{k}}^{2}.
\]
By the two properties of a kept direction,
$g_{k}^{\top}d_{k}\le-\sigma\norm{g_{k}}^{1+p}$ and
$\norm{d_{k}}^{2}\le\kappa^{2}\norm{g_{k}}^{2q}$, so
\[
f(x_{k}) - f(x_{k+1})
\;\ge\; \alpha_{k}\,\sigma\norm{g_{k}}^{1+p}
       - \frac{\Lone}{2}\,\alpha_{k}^{2}\,\kappa^{2}\norm{g_{k}}^{2q}.
\]
Substituting $\alpha_{k} = \tfrac{\sigma}{\Lone\kappa^{2}}\norm{g_{k}}^{1+p-2q}$:
the first term equals
$\tfrac{\sigma^{2}}{\Lone\kappa^{2}}\norm{g_{k}}^{(1+p-2q)+(1+p)}
= \tfrac{\sigma^{2}}{\Lone\kappa^{2}}\norm{g_{k}}^{2(1+p-q)}$,
and the second term equals
$\tfrac{\Lone}{2}\cdot\tfrac{\sigma^{2}}{\Lone^{2}\kappa^{4}}
\norm{g_{k}}^{2(1+p-2q)}\cdot\kappa^{2}\norm{g_{k}}^{2q}
= \tfrac{\sigma^{2}}{2\Lone\kappa^{2}}\norm{g_{k}}^{2(1+p-q)}$.
Subtracting,
\[
f(x_{k}) - f(x_{k+1})
\;\ge\; \frac{\sigma^{2}}{\Lone\kappa^{2}}
        \Bigl(1-\tfrac12\Bigr)\norm{g_{k}}^{2(1+p-q)}
\;=\; \frac{\sigma^{2}}{2\Lone\kappa^{2}}\,\norm{g_{k}}^{2(1+p-q)} .
\qedhere
\]
\end{proof}

\begin{remark}[Optimality of $\alpha^{\ast}_{\setN}$]
The right-hand side of the intermediate bound,
$\alpha\,\sigma\norm{g}^{1+p} - \tfrac{\Lone\kappa^{2}}{2}\alpha^{2}\norm{g}^{2q}$,
is a downward parabola in $\alpha$, maximized exactly at
$\alpha = \tfrac{\sigma}{\Lone\kappa^{2}}\norm{g}^{1+p-2q} = \alpha^{\ast}_{\setN}$.
\end{remark}

\begin{lemma}[Decrease on restarted iterations]\label{lem:decreaseR}
For every restarted iteration $k \in \setR$ (so that $d_{k} = -g_{k}$ and
$\alpha_{k} = 1/\Lone$),
\begin{equation}\label{eq:decreaseR}
f(x_{k}) - f(x_{k+1}) \;\geq\; \cR \, \norm{g_{k}}^{2},
\qquad
\cR = \frac{1}{2\Lone} .
\end{equation}
\end{lemma}

\begin{proof}
Apply \Cref{lem:descent} with $y - x = -\tfrac{1}{\Lone}g_{k}$:
\[
f(x_{k+1})
\;\le\; f(x_{k}) - \frac{1}{\Lone}\norm{g_{k}}^{2}
        + \frac{\Lone}{2}\cdot\frac{1}{\Lone^{2}}\norm{g_{k}}^{2}
\;=\; f(x_{k}) - \frac{1}{2\Lone}\norm{g_{k}}^{2}.
\]
Rearranging gives \eqref{eq:decreaseR}. As in the previous remark, $\alpha = 1/\Lone$
maximizes $\alpha\norm{g}^{2}-\tfrac{\Lone}{2}\alpha^{2}\norm{g}^{2}$, so the
restarted stepsize is also decrease-optimal for its regime.
\end{proof}

% ============================================================
\section{Step II --- Improve or Localize}\label{sec:improve}

\paragraph{Rationale.}
The goal of this section is to bound the total squared displacement
$\sum_{k} \norm{x_{k+1} - x_{k}}^{2}$ by the total function decrease, uniformly over
restarted and non-restarted iterations. We treat the two cases separately and find the
\emph{same} per-step bound, which we then telescope. The resulting inequality admits
the two readings of Step~II in \Cref{sec:roadmap} (improve / localize); the localize
reading is the one consumed by the escape proof.

\subsection{Non-restarted iterations ($k \in \setN$)}\label{sub:improveN}
For a non-restarted iteration the restart condition guarantees
$\norm{d_{k}} \leq \kappa \norm{g_{k}}^{q}$, so
\begin{align}
\norm{x_{k+1} - x_{k}}^{2}
  &= \alpha_{k}^{2} \norm{d_{k}}^{2}
   \;\leq\; \alpha_{k}^{2}\, \kappa^{2} \norm{g_{k}}^{2q} \notag\\
  &= \frac{\sigma^{2}}{\Lone^{2}\kappa^{4}}\,
     \norm{g_{k}}^{\,2(1+p-2q)}\, \kappa^{2}\, \norm{g_{k}}^{2q}
   \;=\; \left(\frac{\sigma}{\Lone\kappa}\right)^{2}
        \norm{g_{k}}^{\,2(1+p-q)} ,
\label{eq:disp-N}
\end{align}
where we used $\alpha_{k} = \tfrac{\sigma}{\Lone\kappa^{2}}\norm{g_{k}}^{1+p-2q}$ and
combined the exponents $2(1+p-2q) + 2q = 2(1+p-q)$.

From \Cref{lem:decreaseN}, $\cN \norm{g_{k}}^{2(1+p-q)} \leq f(x_{k}) - f(x_{k+1})$,
that is,
\begin{equation}\label{eq:grad-bound-N}
\norm{g_{k}}^{\,2(1+p-q)}
  \;\leq\; \frac{2\Lone\kappa^{2}}{\sigma^{2}}\,
           \big(f(x_{k}) - f(x_{k+1})\big) .
\end{equation}
Substituting \eqref{eq:grad-bound-N} into \eqref{eq:disp-N} yields
\begin{equation}\label{eq:disp-N-final}
\norm{x_{k+1} - x_{k}}^{2}
  \;\leq\; \frac{\sigma^{2}}{\Lone^{2}\kappa^{2}}
           \cdot \frac{2\Lone\kappa^{2}}{\sigma^{2}}\,
           \big(f(x_{k}) - f(x_{k+1})\big)
  \;=\; \frac{2}{\Lone}\,\big(f(x_{k}) - f(x_{k+1})\big) .
\end{equation}

\subsection{Restarted iterations ($k \in \setR$)}\label{sub:improveR}
For a restarted iteration $d_{k} = -g_{k}$ and $\alpha_{k} = 1/\Lone$, hence
\begin{equation}\label{eq:disp-R}
\norm{x_{k+1} - x_{k}}^{2}
  = \alpha_{k}^{2}\norm{d_{k}}^{2}
  = \alpha_{k}^{2}\norm{g_{k}}^{2}
  = \frac{1}{\Lone^{2}}\norm{g_{k}}^{2} .
\end{equation}
From \Cref{lem:decreaseR}, $\norm{g_{k}}^{2} \leq \tfrac{1}{\cR}(f(x_{k}) - f(x_{k+1}))
= 2\Lone\,(f(x_{k}) - f(x_{k+1}))$, so
\begin{equation}\label{eq:disp-R-final}
\norm{x_{k+1} - x_{k}}^{2}
  \;\leq\; \frac{1}{\Lone^{2}} \cdot 2\Lone\,
           \big(f(x_{k}) - f(x_{k+1})\big)
  \;=\; \frac{2}{\Lone}\,\big(f(x_{k}) - f(x_{k+1})\big) .
\end{equation}

\subsection{The uniform improve-or-localize lemma}\label{sub:improve-uniform}

\begin{lemma}[Improve or localize]\label{lem:improve}
For every iteration $k$ taken with the optimal stepsize $\alpha^{\ast}$,
\begin{equation}\label{eq:improve-step}
\norm{x_{k+1} - x_{k}}^{2}
  \;\leq\; \frac{2}{\Lone}\,\big(f(x_{k}) - f(x_{k+1})\big) .
\end{equation}
In particular each such step decreases $f$. Consequently, telescoping over any window
of consecutive such steps $k = 0, \dots, T-1$,
\begin{equation}\label{eq:improve-sum}
\sum_{k=0}^{T-1} \norm{x_{k+1} - x_{k}}^{2}
  \;\leq\; \frac{2}{\Lone}\,\big(f(x_{0}) - f(x_{T})\big)
  \;\leq\; \frac{2}{\Lone}\,\big(f(x_{0}) - \fstar\big).
\end{equation}
\end{lemma}

\begin{proof}
The per-step bound \eqref{eq:improve-step} is \eqref{eq:disp-N-final} for
$k \in \setN$ and \eqref{eq:disp-R-final} for $k \in \setR$; the two regimes give the
same constant. The left-hand side of \eqref{eq:improve-step} is nonnegative, so each
per-step decrease $f(x_{k})-f(x_{k+1})$ is nonnegative. Summing
\eqref{eq:improve-step} over $k = 0,\dots,T-1$ and telescoping the right-hand side
gives \eqref{eq:improve-sum}.
\end{proof}

\begin{remark}[Interpretation: a dichotomy]\label{rem:improve}
\Cref{lem:improve} is a dichotomy. At any moment, either the function value has
decreased by a meaningful amount (\emph{improve}), or the iterates did not move much
and we are stuck near where we started (\emph{localize}). The aim of
\Cref{sec:escape} is to upgrade this dichotomy near a strict saddle into the
statement: \emph{at a strict saddle, we always improve}. If we can prove that, then
strict saddles are not an obstruction, and a failure to improve at some iteration can
only mean we are at a (local) minimizer.
\end{remark}

\begin{remark}[Flagged gap: the in-window stepsize is a third case]
\label{rem:window-steps-gap}
% NOTE (flagged correction/gap, do not resolve silently):
\Cref{lem:improve} covers exactly two iteration types: non-restarted steps with
$\alpha = \alpha^{\ast}_{\setN}$, and restarted steps with $\alpha = 1/\Lone$.
Inside the escape window, however, \Cref{algo:PNCG} takes $\alpha_{k} = 1/\Lone$
with \emph{NCG directions} --- a third combination, covered by neither case. For
this combination the per-step decrease can genuinely fail: the descent lemma gives
\[
f(x_{k}) - f(x_{k+1})
\;\ge\; \frac{\sigma}{\Lone}\norm{g_{k}}^{1+p}
        - \frac{\kappa^{2}}{2\Lone}\norm{g_{k}}^{2q},
\]
whose sign depends on $\norm{g_{k}}$. Since \Cref{lem:localization} below applies
\Cref{lem:improve} to in-window steps, this is a gap in the present version of the
argument.

\emph{Candidate resolution:} Set $p = q = 1$, then we can remove the constraint $\alpha_k = 1/\Lone$,
and we still have the same results with some $\alpha_{\min} \in \{\alpha^{\ast}_{\mathcal{N}}, 1/\Lone\}$.
\end{remark}

% ============================================================
\section{Near-Saddle Analysis and Escape}\label{sec:escape}

The improve-or-localize lemma (\Cref{lem:improve}) is generic and
curvature-blind: by itself, it does not exclude the \emph{localize} case at a
strict saddle. To exclude it, we study the dynamics of $u_{k} := x_{k}-\snap$
\emph{centered at the snapshot}, using the second-order Taylor expansion of $f$
around $\snap$. The expansion is accurate in a small ball around $\snap$ ---
precisely the localize regime.

% =============================================================================
\subsection{Step III --- Linearized dynamics around the snapshot}
\label{sub:linearization}
% =============================================================================

\paragraph{Rationale.}
We linearize in order to perform a spectral analysis: get a linear system, compute
its eigenvalues, and show that the system has an eigenvalue strictly larger than
$1$. The element of study is $u_{k} = x_{k}-\snap$, and specifically its projection
$c_{k} = \inner{u_{k}}{\vmin}$ on the negative-curvature direction: if the
escape eigenvalue $\mu_{+}$ exceeds $1$, then $c_{k}$ is amplified geometrically and
the iterate is repelled from the stationary point along $\vmin$. Everything after
this subsection exists to control the non-linear corrections to this picture.

Throughout this section, we condition on a perturbation trigger at iterate
$\snap$, i.e.\
\[
\norm{g_{\snap}} \le \epsg, \qquad \lmin\!\bigl(\hess f(\snap)\bigr) \le -\epsH,
\qquad g_{\snap} := \grad f(\snap).
\]
We center the analysis at $\snap$. Set
\[
H \;:=\; \hess f(\snap),
\]
and let $(\lmin,\vmin)$ denote the eigenpair of $H$ at its most negative
eigenvalue; by the trigger, $|\lmin|\ge\epsH$. Define
\begin{equation}\label{eq:near-saddle-quants}
u_{k} \;:=\; x_{k}-\snap,
\qquad
c_{k} \;:=\; \inner{u_{k}}{\vmin},
\qquad
\eta_{k} \;:=\; \inner{d_{k-1}}{\vmin},
\end{equation}
with the convention $d_{-1}:=0$ (so $\eta_{0}=0$, encoding the
post-perturbation momentum reset $d_{0}=-g_{0}$). Inside the window the stepsize is
constant, $\alpha = 1/\Lone$ (cf.\ \Cref{algo:PNCG}; but see
\Cref{rem:window-steps-gap}).

\paragraph{Quadratic model at $\snap$.}
The second-order Taylor expansion of $f$ at $\snap$ is
\begin{equation}\label{eq:quadmodel}
\tilde f(x) \;=\; f(\snap) + g_{\snap}^{\top}(x-\snap)
   + \tfrac{1}{2}(x-\snap)^{\top}H(x-\snap),
\qquad
\grad\tilde f(x) \;=\; g_{\snap} + H(x-\snap).
\end{equation}
The full gradient differs from the model by the error
\begin{equation}\label{eq:tildee-def}
\tilde e(x) \;:=\; \grad f(x) - g_{\snap} - H(x-\snap),
\end{equation}
bounded in \Cref{lem:graderror} by $\norm{\tilde e(x)}\le(\Ltwo/2)\norm{x-\snap}^{2}$.
We write $\tilde e_{k} := \tilde e(x_{k})$. Unlike the expansion at a critical
point, the constant term $g_{\snap}$ does \emph{not} vanish: the snapshot has
small ($\le\epsg$) but nonzero gradient.

\paragraph{Linearized update.}
The NCG iteration $x_{k+1}=x_{k}+\alpha_{k} d_{k}$ with
$d_{k}=-g_{k}+\beta_{k}d_{k-1}$ and in-window stepsizes
$\alpha_{k}\in\{\amin,1/\Lone\}$ (\Cref{ass:pq}), combined with the
decomposition $g_{k}=g_{\snap}+Hu_{k}+\tilde e_{k}$, gives
\begin{align}
d_{k} &\;=\; -H u_{k} + \beta_{k}d_{k-1} - g_{\snap} - \tilde e_{k},
\label{eq:dk-snap}\\[2pt]
u_{k+1} &\;=\; (I-\alpha_k H)\,u_{k} + \alpha_k\beta_{k}\,d_{k-1}
              - \alpha_k\,g_{\snap} - \alpha_k\,\tilde e_{k}.
\label{eq:uk-snap}
\end{align}
Compared with a recurrence centered at an exact critical point, the new
feature is the constant forcing $-\alpha_k\,g_{\snap}$.

\paragraph{State-space form.}
Stacking $z_{k}:=(u_{k},\,d_{k-1})$,
\begin{equation}\label{eq:state-space}
z_{k+1} \;=\; M_{k}\,z_{k} \;-\; b \;-\; r_{k},
\end{equation}
where
\begin{equation}\label{eq:Mblock}
M_{k} \;=\; \begin{pmatrix}
   I-\alpha H & \alpha\beta_{k}I \\ -H & \beta_{k}I
\end{pmatrix},
\qquad
b \;=\; \begin{pmatrix} \alpha\,g_{\snap} \\ g_{\snap}\end{pmatrix},
\qquad
r_{k} \;=\; \begin{pmatrix} \alpha\,\tilde e_{k} \\ \tilde e_{k}\end{pmatrix}.
\end{equation}
The additive terms --- the constant forcing $b$ from $g_{\snap}$ and the
model-error residual $r_{k}$ --- are exactly the corrections that the later steps
will control.

\paragraph{Diagonalization.}
Since $H$ is symmetric, it diagonalizes in an orthonormal eigenbasis
$\{v_{i}\}_{i=1}^{d}$ with eigenvalues $\lambda_{1}\le\cdots\le\lambda_{d}$
(so $\lambda_{1}=\lmin$ and $v_{1}=\vmin$). Projecting~\eqref{eq:state-space}
onto each $v_{i}$ decouples $M_{k}$ into $d$ independent $2\times2$ blocks
\begin{equation}\label{eq:Miblock}
M_{k,i} \;=\; \begin{pmatrix}
   1-\alpha\lambda_{i} & \alpha\beta_{k} \\ -\lambda_{i} & \beta_{k}
\end{pmatrix}
\end{equation}
acting on the scalar pair $(\inner{u_{k}}{v_{i}},\,\inner{d_{k-1}}{v_{i}})$.
The next lemma shows that, along the negative-curvature direction,
$M_{k,i}$ has an eigenvalue strictly larger than $1$.

\begin{lemma}[Escape eigenvalue along negative curvature]\label{lem:escapeeig}
Fix a step $k$ and an eigenpair $(\lambda_{i},v_{i})$ of $H$ with
$\lambda_{i}<0$, and suppose $\beta_{k}\ge 0$ (as holds for Fletcher-Reeves). The matrix $M_{k,i}$
in~\eqref{eq:Miblock} has a real eigenvalue
\begin{equation}\label{eq:mu-plus}
\mu_{+} \;\ge\; 1 + \alpha\,|\lambda_{i}| \;>\; 1,
\end{equation}
with equality $\mu_{+}=1+\alpha|\lambda_{i}|$ iff $\beta_{k}=0$
(the gradient-descent case).
\end{lemma}

\begin{proof}
The block $M_{k,i}$ has $\operatorname{trace}M_{k,i} = (1-\alpha\lambda_{i})+\beta_{k}$
and
\[
\det M_{k,i}
   \;=\; (1-\alpha\lambda_{i})\,\beta_{k} - (\alpha\beta_{k})(-\lambda_{i})
   \;=\; \beta_{k}.
\]
Its characteristic polynomial is therefore
$p(\mu) = \mu^{2} - (1-\alpha\lambda_{i}+\beta_{k})\mu + \beta_{k}$,
an upward parabola. Evaluating at $\mu=1$,
\[
p(1) \;=\; 1 - (1-\alpha\lambda_{i}+\beta_{k}) + \beta_{k}
     \;=\; \alpha\,\lambda_{i} \;<\; 0,
\]
since $\alpha>0$ and $\lambda_{i}<0$. An upward parabola negative at $\mu=1$
has two distinct real roots straddling $1$, so the larger root exceeds $1$.
For the sharper bound, set $w:=\alpha|\lambda_{i}|=-\alpha\lambda_{i}>0$, so
that $p(\mu)=\mu^{2}-(1+w+\beta_{k})\mu+\beta_{k}$ and
\[
p(1+w) \;=\; (1+w)^{2}-(1+w+\beta_{k})(1+w)+\beta_{k}
        \;=\; -w\,\beta_{k} \;\le\; 0.
\]
Hence $1+w$ lies between the two roots (or equals the larger one when
$\beta_{k}=0$), so the larger root satisfies
$\mu_{+}\ge 1+w$, with equality iff $\beta_{k}=0$. Explicitly,
$\mu_{\pm}=\tfrac{1}{2}\bigl[(1+w+\beta_{k}) \pm
\sqrt{(1+w+\beta_{k})^{2}-4\beta_{k}}\,\bigr]$.
\end{proof}

\paragraph{Projected recurrence on $c_{k}$.}
Projecting~\eqref{eq:uk-snap} on $\vmin$ and using $H\vmin=\lmin\vmin$,
\begin{equation}\label{eq:c-recurrence}
c_{k+1}
\;=\; (1-\alpha_{k}\lmin)\,c_{k} + \alpha_{k}\,\beta_{k}\,\eta_{k}
      - \alpha_{k}\,\inner{g_{\snap}}{\vmin}
      - \alpha_{k}\,\inner{\tilde e_{k}}{\vmin}.
\end{equation}
Since $-\lmin\ge\epsH$, we have $1-\alpha_{k}\lmin\ge1+\alpha_{k}\epsH$. The
reverse triangle inequality applied to~\eqref{eq:c-recurrence} yields
\begin{equation}\label{eq:c-lower}
|c_{k+1}|
\;\ge\; (1+\alpha_{k}\epsH)\,|c_{k}|
\;-\; \alpha_{k}\,\beta_{k}\,|\eta_{k}|
\;-\; \alpha_{k}\,\epsg
\;-\; \alpha_{k}\,\norm{\tilde e_{k}}.
\end{equation}
where we used $|\inner{g_{\snap}}{\vmin}|\le\norm{g_{\snap}}\le\epsg$ and
$|\inner{\tilde e_{k}}{\vmin}|\le\norm{\tilde e_{k}}$.
The recurrence~\eqref{eq:c-lower} contains three perturbing terms: the
\emph{momentum residual} $\alpha\beta_{k}|\eta_{k}|$, the \emph{snapshot
forcing} $\alpha\epsg$, and the \emph{model error} $\alpha\norm{\tilde e_{k}}$.
Each will be controlled separately in the inductive analysis
(\Cref{ssec:induction}).

\begin{remark}[Interpretation, and what remains to be done]\label{rem:escape}
The escape eigenvalue $\mu_{+}>1$ from \Cref{lem:escapeeig} drives geometric
growth of $c_{k}$ under the pure linearization (i.e., $\tilde e\equiv 0$ and
$g_{\snap}=0$): the iterate is repelled from the snapshot along the
negative-curvature direction $\vmin$. In the actual NCG iteration, three
corrections must be handled:
\begin{itemize}[leftmargin=*]
   \item the \emph{model error} $\tilde e_{k}$, controlled by
     \Cref{lem:graderror} (Step~IV);
   \item the \emph{snapshot forcing} $g_{\snap}$, small by the trigger
     condition $\norm{g_{\snap}}\le\epsg$;
   \item the \emph{absence of an initial signal}: the geometric growth
     $|c_{k}|\gtrsim\mu_{+}^{k}|c_{0}|$ is vacuous if $|c_{0}|$ is zero (or
     tiny). The random perturbation $\xi$ guarantees a nontrivial initial
     $|c_{0}| = |\inner{\xi}{\vmin}|$ via the anti-concentration estimate of
     \Cref{lem:anticoncentration} (Step~V).
\end{itemize}
The momentum residual additionally requires an \emph{upper bound on
$\beta_{k}$}, which for Fletcher--Reeves requires a gradient \emph{lower}
bound (Step~VI). The remainder of \Cref{sec:escape} formalizes each of these
corrections, and Step~VII assembles them into the induction.
\end{remark}

% =============================================================================
\subsection{Step IV --- Quadratic model error}\label{sub:graderror-section}
% =============================================================================

\paragraph{Rationale.}
Straightforward step: bound $\norm{\tilde e_{k}}$ above. This is the error between
the real gradient and the quadratically approximated one; if it could exceed the
geometric push $\alpha\epsH|c_{k}|$, the escape signal could be cancelled. The
bound is quadratic in the distance to the snapshot, which is why the analysis
must keep the iterates confined (invariant~(i) of Step~VII).

\begin{lemma}[Quadratic model error at the snapshot]\label{lem:graderror}
Under \Cref{ass:hesslip}, for every $x\in\R^{d}$,
\begin{equation}\label{eq:graderror-snap}
\grad f(x) \;=\; g_{\snap} + H(x-\snap) + \tilde e(x),
\qquad
\norm{\tilde e(x)} \;\le\; \frac{\Ltwo}{2}\,\norm{x-\snap}^{2}.
\end{equation}
\end{lemma}

\begin{proof}
Let $x_{t}:=\snap+t(x-\snap)$ for $t\in[0,1]$, so $\frac{dx_{t}}{dt}=x-\snap$.
Set $\phi(t):=\grad f(x_{t})$; by the chain rule,
$\phi'(t)=\hess f(x_{t})(x-\snap)$. By the fundamental theorem of calculus,
\[
\grad f(x)-g_{\snap}
\;=\; \grad f(x)-\grad f(\snap)
\;=\; \int_{0}^{1}\phi'(t)\,dt
\;=\; \int_{0}^{1}\hess f(x_{t})(x-\snap)\,dt.
\]
Subtracting $\int_{0}^{1}H(x-\snap)\,dt = H(x-\snap)$,
\[
\tilde e(x)
\;=\; \int_{0}^{1}\bigl(\hess f(x_{t})-H\bigr)(x-\snap)\,dt.
\]
Taking norms and using \Cref{ass:hesslip} together with $x_{t}-\snap=t(x-\snap)$,
\[
\norm{\tilde e(x)}
\;\le\; \int_{0}^{1}\norm{\hess f(x_{t})-\hess f(\snap)}\,\norm{x-\snap}\,dt
\;\le\; \int_{0}^{1}\Ltwo\,t\,\norm{x-\snap}^{2}\,dt
\;=\; \frac{\Ltwo}{2}\,\norm{x-\snap}^{2}.
\qedhere
\]
\end{proof}

% =============================================================================
\subsection{Setting for the probabilistic escape}\label{sub:setting}
% =============================================================================

The post-perturbation iterate is
\[
x_{0} \;:=\; \snap + \xi,
\qquad
\xi \sim \mathcal{U}\!\bigl(\mathcal{B}^{d}(0,\Rstar)\bigr),
\]
so that $u_{0}=\xi$ and $\norm{u_{0}}\le \Rstar$ deterministically. The escape window runs PR-NCG for $T=\Tp$ steps with the algorithm's ordinary
stepsize rule, so that $\alpha_{k}\in\{\amin,\,1/\Lone\}$ at every window step
(\Cref{ass:pq}); in particular $\alpha_{k}\ge\amin$ and
$\alpha_{k}\Lone\le 1$ throughout.

The three distinguished constants of the analysis are the dimensional factor, the
\emph{optimal perturbation radius}, and the \emph{anti-concentration threshold}:
\begin{equation}\label{eq:Rstar-snap}
a \;:=\; \sqrt{\frac{\pi}{2(d+2)}},
\qquad
\Rstar \;:=\; \frac{a\,\epsH\,\delta}{\Ltwo},
\qquad
\sdel \;:=\; a\,\delta\,\Rstar
\;=\; \frac{a^{2}\,\epsH\,\delta^{2}}{\Ltwo}
\;=\; \frac{\pi\,\epsH\,\delta^{2}}{2\,\Ltwo\,(d+2)}.
\end{equation}
The choice $R = \Rstar$ is justified by the optimization carried out in the
proof of \Cref{lem:gfloor0} (Step~VI). Direct computation from
\eqref{eq:Rstar-snap} gives the exact identities
\begin{equation}\label{eq:identities}
\frac{\Rstar}{\sdel} \;=\; \frac{1}{a\,\delta},
\qquad
\Ltwo\,(\Rstar)^{2} \;=\; \epsH\,\sdel ,
\end{equation}
which will be used repeatedly in Step~IX (the second identity is the
\emph{structural identity} of \Cref{lem:structural}).

% =============================================================================
\subsection{Step V --- Anti-concentration of the perturbation}\label{sub:anticonc}
% =============================================================================

\paragraph{Rationale.}
Straightforward step: show that the perturbation has a non-negligible component
along the negative-curvature direction, with high probability. This is the only
probabilistic ingredient in the escape analysis; everything downstream is
deterministic conditional on the good event $\mathcal{G}_{0}$.

\begin{lemma}[Anti-concentration]\label{lem:anticoncentration}
Let $\xi\sim\mathcal{U}(\mathcal{B}^{d}(0,R))$ and let $v\in\R^{d}$ be a unit
vector. For every $s\in(0,R]$,
\begin{equation}\label{eq:anticonc-prob}
\Pr\bigl(|\inner{\xi}{v}|\le s\bigr)
\;\le\; \frac{s}{R}\,\sqrt{\frac{2(d+2)}{\pi}}.
\end{equation}
In particular, with $R=\Rstar$ as in~\eqref{eq:Rstar-snap} and $s=\sdel$, the event
\[
\mathcal{G}_{0} \;:=\; \bigl\{|\inner{\xi}{\vmin}|\ge\sdel\bigr\}
\qquad\text{satisfies}\qquad
\Pr(\mathcal{G}_{0}) \;\ge\; 1-\delta.
\]
\end{lemma}

\begin{proof}
\emph{Step 1: rotational invariance reduction.}
The Lebesgue measure on $\mathcal{B}^{d}(0,R)$ is invariant under any
orthogonal $Q\in O(d)$, hence $Q\xi\stackrel{d}{=}\xi$. Choose $Q$ with
$Qe_{1}=v$ (complete $v$ to an orthonormal basis). Then
\[
\inner{\xi}{v} \;=\; \inner{\xi}{Qe_{1}} \;=\; \inner{Q^{\top}\xi}{e_{1}},
\qquad Q^{\top}\xi\stackrel{d}{=}\xi,
\]
so the distribution of $\inner{\xi}{v}$ depends only on $\norm{v}=1$, not on
the direction of $v$. Take $v=e_{1}$ and set $X:=\xi_{1}$.

\emph{Step 2: marginal density of $X$.}
Write $\xi=(\xi_{1},\xi_{\perp})$ with $\xi_{\perp}\in\R^{d-1}$. By Fubini's
theorem on the uniform density,
\[
f_{X}(x)
\;=\; \int_{\R^{d-1}}\frac{\indic{x^{2}+\norm{\xi_{\perp}}^{2}\le R^{2}}}{V_{d}R^{d}}\,d\xi_{\perp},
\]
with $V_{k}:=\pi^{k/2}/\Gamma(k/2+1)$ the unit-ball volume in $\R^{k}$. For
$|x|\le R$, the slice is the $(d-1)$-ball of radius $\sqrt{R^{2}-x^{2}}$, of
measure $V_{d-1}(R^{2}-x^{2})^{(d-1)/2}$, so
\begin{equation}\label{eq:marginal-density}
f_{X}(x) \;=\; \frac{V_{d-1}}{V_{d}R^{d}}\,(R^{2}-x^{2})^{(d-1)/2}\,\indic{|x|\le R}.
\end{equation}

\emph{Step 3: density peaks at $0$.}
For $|x|\le R$, $(R^{2}-x^{2})^{(d-1)/2}$ is decreasing in $|x|$, hence
\[
f_{X}(x)\;\le\; f_{X}(0) \;=\; \frac{V_{d-1}}{V_{d}\,R}.
\]

\emph{Step 4: volume ratio via Gautschi's inequality.}
By definition,
\[
\frac{V_{d-1}}{V_{d}}
\;=\; \frac{1}{\sqrt{\pi}}\cdot\frac{\Gamma(d/2+1)}{\Gamma((d+1)/2)}.
\]
Recall Gautschi's inequality: for $x>0$ and $0<s<1$,
$\Gamma(x+1)/\Gamma(x+s)\le(x+1)^{1-s}$. With $x=d/2$, $s=1/2$,
\[
\frac{\Gamma(d/2+1)}{\Gamma((d+1)/2)} \;\le\; \sqrt{(d+2)/2},
\qquad\text{so}\qquad
\frac{V_{d-1}}{V_{d}} \;\le\; \sqrt{\frac{d+2}{2\pi}}.
\]

\emph{Step 5: integrate.}
Combining Steps 3 and 4, $f_{X}(0)\le \frac{1}{R}\sqrt{(d+2)/(2\pi)}$, so
\[
\Pr(|X|\le s) \;=\; \int_{-s}^{s}f_{X}(x)\,dx
\;\le\; 2s\,f_{X}(0)
\;\le\; \frac{s}{R}\,\sqrt{\frac{2(d+2)}{\pi}},
\]
which is~\eqref{eq:anticonc-prob}.

\emph{Step 6: high-probability event.}
With $R=\Rstar$ and $s=\sdel=a\delta \Rstar$, recalling
$a = \sqrt{\pi/(2(d+2))}$,
\[
\Pr\bigl(|\inner{\xi}{\vmin}|\le\sdel\bigr)
\;\le\; \frac{\sdel}{\Rstar}\sqrt{\frac{2(d+2)}{\pi}}
\;=\; a\,\delta\cdot\frac{1}{a}
\;=\; \delta,
\]
so $\Pr(\mathcal{G}_{0})=1-\Pr(|\inner{\xi}{\vmin}|\le\sdel)\ge 1-\delta$.
\end{proof}

\begin{remark}[Reading \Cref{lem:anticoncentration}]
This lemma is the only probabilistic ingredient in the escape analysis.
It guarantees that, with probability at least $1-\delta$, the random
perturbation has a non-negligible component along the negative-curvature
direction: $|\inner{\xi}{\vmin}|\ge\sdel$. A~priori the uniform
perturbation could land almost orthogonally to $\vmin$, giving
$\inner{\xi}{\vmin}\approx 0$ and no usable escape signal, in which case
the geometric growth from \Cref{lem:escapeeig} would have nothing to
amplify. \Cref{lem:anticoncentration} excludes this bad case with high
probability; everything that follows is deterministic conditional on the
event $\mathcal{G}_{0}$.
\end{remark}

% =============================================================================
\subsection{Step VI --- Post-perturbation gradient lower bound}\label{sub:gfloor}
% =============================================================================

\paragraph{Rationale (Fletcher--Reeves specific).}
This step exists to control the momentum term. Since
$\beta^{\mathrm{FR}}_{k} = \norm{g_{k}}^{2}/\norm{g_{k-1}}^{2}$, bounding
$\beta^{\mathrm{FR}}$ requires bounding $\norm{g_{k}}$ \emph{above} and
$\norm{g_{k-1}}$ \emph{below}:
$\beta^{\mathrm{FR}} \le \gup^{2}/\glow^{2}$. The upper bound is easy
(Lipschitz, cf.~\eqref{eq:barg-def}); the lower bound is the hard part, because
Lipschitz continuity only ever produces upper bounds, and the window sits in a
region where the gradient can in principle vanish. Any other $\beta$ formula
would mean redoing this step; heavy ball or Nesterov do not need it, because
their momentum coefficient is constant. This subsection establishes the lower
bound at time $0$ (and, as a by-product, fixes the optimal perturbation radius
$\Rstar$); propagating the lower bound across the whole window is the job of the
induction (Step~VII).

\begin{lemma}[Post-perturbation gradient lower bound]\label{lem:gfloor0}
Suppose
\begin{equation}\label{eq:smallness-snap}
\epsg \;<\; \tfrac{1}{2}\,\epsH\,\sdel.
\end{equation}
Then on the event $\mathcal{G}_{0}$ of \Cref{lem:anticoncentration},
\[
\norm{g_{0}} \;\ge\; \glow_{0},
\qquad
\glow_{0} \;:=\; \tfrac{1}{2}\,\epsH\,\sdel - \epsg \;>\; 0.
\]
\end{lemma}

\begin{proof}
\emph{Step 1: quadratic-model decomposition.}
By \Cref{lem:graderror} applied at $x_{0}=\snap+\xi$ and $u_{0}=\xi$,
\[
g_{0} \;=\; g_{\snap} + H\,\xi + \tilde e_{0},
\qquad
\norm{\tilde e_{0}} \;\le\; \tfrac{\Ltwo}{2}\,\norm{\xi}^{2}
\;\le\; \tfrac{\Ltwo}{2}\,R^{2},
\]
where $R$ is the perturbation radius (to be optimized in Step~5 below).

\emph{Step 2: project on $\vmin$.}
Using $H\vmin=\lmin\vmin$,
\begin{equation}\label{eq:g0proj}
\inner{g_{0}}{\vmin}
\;=\; \underbrace{\inner{g_{\snap}}{\vmin}}_{=:C}
\;+\; \underbrace{\lmin\,\inner{\xi}{\vmin}}_{=:P}
\;+\; \underbrace{\inner{\tilde e_{0}}{\vmin}}_{=:E}.
\end{equation}

\emph{Step 3: bound the three terms.}
On $\mathcal{G}_{0}$, $|\inner{\xi}{\vmin}|\ge\sdel$; combined with
$|\lmin|\ge\epsH$,
\[
|P| \;\ge\; \epsH\,\sdel,
\qquad
|C| \;\le\; \norm{g_{\snap}} \;\le\; \epsg,
\qquad
|E| \;\le\; \norm{\tilde e_{0}} \;\le\; \tfrac{\Ltwo}{2}\,R^{2}.
\]

\emph{Step 4: combine via reverse triangle inequality.}
From $\inner{g_{0}}{\vmin}=C+P+E$ and $|a+b+c|\ge|a|-|b|-|c|$,
\begin{equation}\label{eq:g0-raw}
\norm{g_{0}}
\;\ge\; |\inner{g_{0}}{\vmin}|
\;\ge\; |P|-|C|-|E|
\;\ge\; \epsH\,\sdel - \epsg - \tfrac{\Ltwo}{2}\,R^{2}.
\end{equation}
Note $\sdel=a\delta R$ depends on $R$; substituting,
\begin{equation}\label{eq:g0-h}
\norm{g_{0}} \;\ge\; h(R)-\epsg,
\qquad
h(R) \;:=\; a\,\epsH\,\delta\,R - \frac{\Ltwo}{2}\,R^{2}.
\end{equation}
The map $R\mapsto h(R)$ is a downward parabola.

\emph{Step 5: optimize in $R$.}
Setting $h'(R)=0$,
\[
h'(R) \;=\; a\,\epsH\,\delta - \Ltwo\,R
\quad\Longrightarrow\quad
\Rstar \;=\; \frac{a\,\epsH\,\delta}{\Ltwo},
\]
matching~\eqref{eq:Rstar-snap}. Compute $h(\Rstar)$ in two pieces:
\[
(\Rstar)^{2}
\;=\; \frac{a^{2}\,\epsH^{2}\,\delta^{2}}{\Ltwo^{2}},
\qquad
\frac{\Ltwo}{2}(\Rstar)^{2}
\;=\; \frac{a^{2}\,\epsH^{2}\,\delta^{2}}{2\,\Ltwo}
\;=\;\tfrac12\,\epsH\,\sdel,
\qquad
a\,\epsH\,\delta\,\Rstar
\;=\; \frac{a^{2}\,\epsH^{2}\,\delta^{2}}{\Ltwo}
\;=\;\epsH\,\sdel,
\]
where we used $\sdel = a^{2}\epsH\delta^{2}/\Ltwo$
from~\eqref{eq:Rstar-snap}. (Note in passing that the middle computation is the
structural identity $\Ltwo(\Rstar)^{2}=\epsH\sdel$
of~\eqref{eq:identities}.) Summing,
\[
h(\Rstar)
\;=\; \epsH\,\sdel-\tfrac12\,\epsH\,\sdel
\;=\; \tfrac{1}{2}\,\epsH\,\sdel .
\]

\emph{Step 6: conclude.}
Setting $R=\Rstar$ in~\eqref{eq:g0-h},
\[
\norm{g_{0}} \;\ge\; \tfrac{1}{2}\,\epsH\,\sdel - \epsg \;=\; \glow_{0},
\]
strictly positive under the smallness condition~\eqref{eq:smallness-snap}.
\end{proof}

\begin{remark}[Why this lemma matters]
If $\norm{g_{0}}$ could vanish, the algorithm could trigger another perturbation
immediately and the whole escape window would be wasted; more importantly, the FR
coefficient $\beta_{1}$ would be uncontrolled. \Cref{lem:gfloor0} rules this out at
time $0$. The window-wide floor $\glow$ of \eqref{eq:glow-def}, established
inductively in \Cref{lem:induction}, is the time-$k$ analogue; note that
$\glow \le \glow_{0}$, since the window floor pays the model error at the larger
radius $S+\Rstar$ rather than $\Rstar$ (compare \eqref{eq:glow-def} with
$\glow_{0} = \epsH\sdel - \epsg - \tfrac{\Ltwo}{2}(\Rstar)^{2}$, using the
structural identity).
\end{remark}

% =============================================================================
\subsection{Step VII --- Setting up the escape window}\label{ssec:setting}
% =============================================================================

\paragraph{Rationale.}
The single perturbation has placed the iterate at $x_0 = \snap + \xi$, and
the algorithm now runs for $T$ steps. Our goal is to prove that, with
high probability, the function value drops by at least $F$ at some point
within this window. Following the improve-or-localize dichotomy
(\Cref{rem:improve}), there are two cases:
\begin{enumerate}[label=(\alph*)]
\item the function value drops by at least $F$ somewhere in $\{0,\dots,T\}$ ---
this is what we try to prove;
\item the function value never drops by that much, in which case
\Cref{lem:improve} tells us the iterates stay close to $x_{0}$
($\norm{x_{k}-x_{0}}\le S$, $S=\sqrt{2FT/\Lone}$).
\end{enumerate}
We work in case~(b) and derive a contradiction, so we must be in case~(a). Case~(b)
is the bad case, and we formalize it by the localization event below.

\paragraph{The localization event.}
Let $T\in\N$ and $F>0$ be parameters (their values will be fixed in
\Cref{thm:escape-single} and Step~IX). Define the
\emph{localization event}
\begin{equation}\label{eq:LT}
\mathcal{L}_T \;:=\; \Bigl\{\max_{0\le k\le T}\bigl(f(x_0)-f(x_k)\bigr) \;<\; F\Bigr\}.
\end{equation}
The event $\mathcal{L}_T$ is the bad case: the function value has not dropped
by $F$ anywhere in the window. Its complement is exactly the conclusion we
want. So the strategy is to show that $\mathcal{G}_{0}\cap\mathcal{L}_{T}$ is
empty: throughout this section and the next, we work conditional on
$\mathcal{G}_0 \cap \mathcal{L}_T$ and derive a contradiction; combined with
$\Pr(\mathcal{G}_0) \ge 1-\delta$ this gives the result.

\paragraph{Confinement on $\mathcal{L}_T$.}
On $\mathcal{L}_T$ the improve-or-localize lemma (\Cref{lem:improve}) forces
the iterates to stay near $x_0$. Recentering at the snapshot yields the
confinement we need.

\begin{lemma}[Localized confinement around $\snap$]\label{lem:localization}
Set the localization radius
\begin{equation}\label{eq:S-def}
S \;:=\; \sqrt{\frac{2FT}{\Lone}}.
\end{equation}
On the event $\mathcal{L}_T$, the displacement from the snapshot satisfies
\begin{equation}\label{eq:confine}
\norm{u_k} \;=\; \norm{x_k-\snap} \;\le\; S + \Rstar
\qquad\text{for every }k\in\{0,1,\dots,T\}.
\end{equation}
\end{lemma}

\begin{proof}
Fix $k\in\{0,\dots,T\}$; the case $k=0$ is immediate from
$\norm{u_{0}}=\norm{\xi}\le R^{*}$. Writing the displacement as a telescoping
sum, then applying the triangle inequality and Cauchy--Schwarz on the sum of
$k$ terms,
\[
\norm{x_{k}-x_{0}}^{2}
\;=\;\Bigl\|\sum_{j=0}^{k-1}(x_{j+1}-x_{j})\Bigr\|^{2}
\;\le\;\Bigl(\sum_{j=0}^{k-1}\norm{x_{j+1}-x_{j}}\Bigr)^{2}
\;\le\; k\,\sum_{j=0}^{k-1}\norm{x_{j+1}-x_{j}}^{2}.
\]
Every window step is either restarted or a non-restarted step with
$\alpha^{\ast}_{\setN}$ (\Cref{ass:pq}, \Cref{rem:safeguard-changelog}), so
the per-step bound \eqref{eq:improve-step} of \Cref{lem:improve} applies at
every $j$; its right-hand sides are nonnegative and telescope:
\[
\sum_{j=0}^{k-1}\norm{x_{j+1}-x_{j}}^{2}
\;\le\;\frac{2}{\Lone}\bigl(f(x_{0})-f(x_{k})\bigr)
\;<\;\frac{2F}{\Lone}
\qquad\text{on }\mathcal{L}_{T}.
\]
Combining, and using $k\le T$,
$\norm{x_{k}-x_{0}}^{2} < 2kF/\Lone \le 2TF/\Lone = S^{2}$. The triangle
inequality with $\norm{x_{0}-\snap}=\norm{\xi}\le R^{*}$ gives
$\norm{u_{k}}\le\norm{x_{k}-x_{0}}+\norm{x_{0}-\snap}\le S+R^{*}$.
\end{proof}

The two contributions to $S+\Rstar$ are independent: $S$ tracks how far the
iterate may wander inside the localization ball, while $\Rstar$ accounts for
the initial perturbation displacement. This is why the analysis is posed
with respect to $S+\Rstar$ rather than $S$ alone.
(Caveat: the application of \Cref{lem:improve} to in-window steps is subject to
the gap flagged in \Cref{rem:window-steps-gap}.)

\paragraph{Validity of the quadratic model.}
At distance $\rho$ from $\snap$, the negative-curvature signal of $\hess
f(\snap)$ contributes roughly $\epsH\,\rho$ to the gradient in direction
$\vmin$, while the quadratic-model error of \Cref{lem:graderror} grows as
$\tfrac{\Ltwo}{2}\rho^2$. For the signal to dominate the error throughout
the confinement ball, we need $\epsH\rho \ge \tfrac{\Ltwo}{2}\rho^2$ at the
worst point $\rho = S+\Rstar$. This is the validity condition
\begin{equation}\label{eq:S-validity}
S + \Rstar \;\le\; \frac{2\epsH}{\Ltwo}.
\tag{C-valid}
\end{equation}
Beyond this threshold, the Lipschitz-Hessian variation swamps the curvature
signal and the linear model is no longer reliable. It will turn out
(\Cref{rem:cvalid}) that \eqref{eq:S-validity} is \emph{implied} by the
constraint (C-err) of \Cref{lem:induction}, so it does not need to be carried
separately in the constants ledger of Step~IX.

\paragraph{Uniform gradient upper bound.}
The momentum analysis requires bounding the FR coefficient $\beta_k =
\norm{g_k}^2/\norm{g_{k-1}}^2$, which in turn requires both an upper and a
lower bound on gradient norms across the window. The upper bound is
immediate from gradient-Lipschitz: applying \Cref{ass:gradlip} to the pair
$(x_k, \snap)$ and combining with the trigger bound $\norm{\grad
f(\snap)}\le\epsg$ and the confinement \eqref{eq:confine},
\begin{equation}\label{eq:barg-def}
\norm{g_k}
\;\le\; \norm{g_{\snap}} + \Lone\norm{x_{k}-\snap}
\;\le\; \epsg + \Lone\,(S+\Rstar) \;=:\; \gup
\qquad\text{for every }k\in\{0,\dots,T\}\text{ on }\mathcal{L}_T.
\end{equation}
The first term $\epsg$ accounts for the residual gradient at the snapshot;
the second term $\Lone(S+\Rstar)$ accounts for the additional gradient that
develops as the iterate moves through the confinement ball. In the relevant
regime $\epsg \ll \Lone(S+\Rstar)$, the second term dominates.

The complementary lower bound cannot come from Lipschitz alone --- Lipschitz
gives only upper bounds. In fact \emph{no} bound over the confinement ball can
supply it: the ball is centered at a near-critical point, so
$\inf_{x\in\mathcal{B}(\snap,\,S+\Rstar)}\norm{\grad f(x)}$ is essentially zero.
The lower bound is a property of \emph{where the iterates actually go}, not of
the region they are in; it requires that $|c_k|$ stays large and that the model
error stays small, both of which are properties we establish in the induction.
We therefore defer the lower bound to that section, where it appears as the
gradient-floor invariant.

\paragraph{Gradient-floor placeholder.}
Anticipating the induction, define the gradient floor
\begin{equation}\label{eq:glow-def}
\glow \;:=\; \epsH\sdel - \epsg - \tfrac{\Ltwo}{2}(S+\Rstar)^{2}.
\end{equation}
We will show in \Cref{lem:induction} that $\glow > 0$ under the budget-split
constraints, and that $\norm{g_k} \ge \glow$ for every $k$ in the window on
$\mathcal{G}_0 \cap \mathcal{L}_T$. Until then, $\glow$ is a notational
placeholder for the quantity we will eventually establish as a uniform lower
bound. Together, $\gup$ and $\glow$ yield the FR envelope
\begin{equation}\label{eq:bbeta-def}
\beta_{k} \;=\; \frac{\norm{g_{k}}^{2}}{\norm{g_{k-1}}^{2}}
\;\le\; \frac{\gup^{2}}{\glow^{2}} \;=:\; \bbeta ,
\end{equation}
used in the momentum constraint (C-mom) below.

\paragraph{Budget splits for the geometric driver.}
The recurrence~\eqref{eq:c-lower} for $c_k$ has the shape
$|c_{k+1}| \ge (1+\alpha\epsH)|c_k| - (\text{three perturbing terms})$, with
contributions from the snapshot forcing $g_{\snap}$, the quadratic-model
error $\tilde e_k$, and the momentum term $\beta_k\eta_k$. To extract
geometric growth, we will require each of these to be a controlled fraction
of the driving term $\alpha\epsH|c_k|$. Introduce three positive splits
$\theta,\theta',\theta'' \in (0,1)$ with $\theta+\theta'+\theta''<1$, and
define the residual growth rate
\begin{equation}\label{eq:gamma-def}
\gamma \;:=\; 1 - \theta - \theta' - \theta''.
\end{equation}
The values of the splits are \emph{not} fixed here: they are decision
variables of the analysis. Their optimal values are derived exactly in the
complexity section (\Cref{def:canonical}); until then every statement holds
for any admissible choice.
% NOTE (flagged correction): an earlier version fixed the "canonical" splits to
% theta = theta' = theta'' = 1/6 at this point. That choice is superseded by the
% exact optimization of Step IX (Definition \ref{def:canonical}); fixing values
% here would contradict it.

\paragraph{Summary of the setup.}
We have committed to: a localization event $\mathcal{L}_T$ on which we work;
a localization radius $S$ and a confinement bound $S+\Rstar$ around $\snap$;
a quadratic-model validity constraint \eqref{eq:S-validity} (later absorbed
into (C-err)); an upper bound $\gup$ on gradient norms; budget splits
$\theta,\theta',\theta''$; and a placeholder gradient floor $\glow$. The
next subsection states the joint invariants and proves them by induction.

% =============================================================================
\subsection{Step VII --- Joint invariants on the escape window}
\label{ssec:induction}
% =============================================================================

\paragraph{Rationale: why an induction is unavoidable.}
The recurrence \eqref{eq:c-lower},
\[
|c_{k+1}| \;\ge\; (1+\alpha\epsH)\,|c_k| \;-\; \alpha\epsg
\;-\;\alpha\beta_{k}|\eta_{k}|\;-\;\alpha\norm{\tilde e_{k}},
\]
cannot be iterated directly, because the obstacle terms depend on the past:
\begin{itemize}[leftmargin=*]
\item $\norm{\tilde e_{k}} \le \tfrac{\Ltwo}{2}\norm{u_{k}}^{2}$ depends on
$u_{k} = x_{k}-\snap$;
\item $\beta_{k}$ depends on $\norm{g_{k-1}}$, which depends on
$\norm{u_{k-1}}$ and on $|c_{k-1}|$;
\item $|\eta_{k}| = |\inner{d_{k-1}}{\vmin}|$ depends on the previous
directions.
\end{itemize}
The solution is an induction, which lets us maintain a set of properties
simultaneously up to time $k$, and prove them at time $k+1$. From the
recurrence, the three things we need are:
\begin{enumerate}
\item how far the iterate is: $\norm{x_{k}-\snap}\le S+\Rstar$;
\item how big the gradient is: $\norm{g_{k}}\ge\glow$;
\item how big $|c_{k}|$ is: $|c_{k}|\ge(1+\gamma\alpha\epsH)^{k}|c_{0}|$.
\end{enumerate}
It is worth pausing on why no shortcut exists for item~2. The upper envelope
$\gup$ is a supremum over the confinement ball --- a static quantity, valid for
any point of the ball. The lower bound has no static analogue: the infimum of
$\norm{\grad f}$ over a ball centered at a near-critical point is essentially
zero. The floor holds only along the trajectory, and only \emph{because} the
trajectory carries a large escape component:
$\norm{g_{k}}\ge\epsH|c_{k}|-\epsg-\norm{\tilde e_{k}}$. This creates the
circular dependency
\[
\text{floor at } j<k
\;\Longrightarrow\; \beta_{j}\text{ controlled}
\;\Longrightarrow\; \text{growth of } |c_{j}| \text{ up to } k
\;\Longrightarrow\; \text{floor at } k,
\]
which is broken precisely by carrying all three properties together in the
inductive hypothesis: each property at step $k+1$ is purchased with a
\emph{different} property at step $k$. Without~(1) the model error blows up;
without~(2) the FR coefficient $\beta_{k}$ blows up; without~(3) there is no
floor and no escape. They prop each other up.

The three constraints (C-snap), (C-err), (C-mom) below each have the form
``\emph{this perturbing source is bounded by
$\theta_{\bullet}\cdot\epsH\sdel$}''. Combined with the inductive lower
bound $|c_{k}|\ge\sdel$, each constraint absorbs the corresponding
perturbing term into a $\theta_{\bullet}$-fraction of the budget
$\alpha\epsH|c_{k}|$, leaving the residual fraction
$\gamma = 1-\theta-\theta'-\theta''$ as the geometric growth rate.

\begin{lemma}[Window invariants for the escape]\label{lem:induction}
Let $p=q=1$ (\Cref{ass:pq}) and suppose the three budget-split constraints
\begin{align}
\epsg                                &\;\le\; \theta''\,\epsH\,\sdel &&\textnormal{(C-snap)}\label{eq:snap-split}\\
\tfrac{\Ltwo}{2}(S+R^{*})^{2}         &\;\le\; \theta\,\epsH\,\sdel   &&\textnormal{(C-err)}\label{eq:err-split}\\
\bar\beta\,\kappa\,\bar g        &\;\le\; \theta'\,\epsH\,\sdel  &&\textnormal{(C-mom)}\label{eq:mom-split}
\end{align}
are satisfied. Then $\glow > 0$ and, on $\mathcal{G}_{0}\cap\mathcal{L}_{T}$,
for every $k\in\{0,\dots,T\}$:
\begin{enumerate}[label=\textnormal{(\roman*)}]
\item $\norm{u_{k}} \;\le\; S+R^{*}$ \hfill \textnormal{(confinement)},
\item $\norm{g_{k}} \;\ge\; \glow$ \hfill \textnormal{(gradient floor)},
\item $|c_{k}| \;\ge\; \prod_{j=0}^{k-1}\bigl(1+\gamma\,\alpha_{j}\,\epsH\bigr)\,|c_{0}|
\;\ge\; (1+\gamma\,\amin\,\epsH)^{k}\,|c_{0}|$
with $|c_{0}|\ge\sdel$ \hfill \textnormal{(geometric growth)}.
\end{enumerate}
\end{lemma}

\begin{proof}
\proofpara{Positivity of $\glow$.}
By \eqref{eq:glow-def} and the constraints (C-snap), (C-err),
\[
\glow \;=\; \epsH\sdel - \epsg - \tfrac{\Ltwo}{2}(S+\Rstar)^2
\;\ge\; (1 - \theta - \theta'')\,\epsH\,\sdel \;>\; 0,
\]
since $\theta+\theta'' < 1$. The remainder of the proof is by induction on $k$:
the inductive hypothesis at step $k$ is that (i)--(iii) hold at every index
$j\in\{0,\dots,k\}$ (strong form; this is what makes $\beta_{k}$, which
involves $\norm{g_{k-1}}$, controllable).

\proofpara{Base case ($k=0$).}
\emph{Confinement.} $u_0 = \xi$, so $\norm{u_0} = \norm{\xi} \le \Rstar \le S+\Rstar$.

\emph{Growth.} $c_0 = \inner{\xi}{\vmin}$ and the event $\mathcal{G}_0$ is
defined precisely as $|c_0|\ge\sdel$.

\emph{Floor.} Project the model identity~\eqref{eq:graderror-snap} at $x_0$
on $\vmin$ and apply the reverse triangle inequality with $|\lmin|\ge\epsH$,
the snapshot trigger $\norm{g_{\snap}}\le\epsg$, and the growth bound just
established:
\[
\norm{g_0}
\;\ge\; |\lmin||c_0| - \norm{g_{\snap}} - \norm{\tilde e_0}
\;\ge\; \epsH\sdel - \epsg - \tfrac{\Ltwo}{2}(\Rstar)^2
\;\ge\; \glow,
\]
where the last step uses $(\Rstar)^2\le(S+\Rstar)^2$ in the definition of $\glow$.

\proofpara{Inductive step.}
Assume invariants (i), (ii), (iii) hold at every index $j \le k$ for some
$0 \le k < T$. We establish them at step $k+1$ in the order: (i) confinement,
then growth (iii), then floor (ii). Each sub-claim uses specific tools, which we
identify line-by-line.

\subparagraph{Confinement at $k+1$.}
\emph{Claim.} $\|u_{k+1}\| \le S + \Rstar$.

\emph{Why this is direct from \Cref{lem:localization}, not from the recurrence.}
\Cref{lem:localization} establishes $\|u_j\| \le S + \Rstar$ \emph{uniformly}
for every $j \in \{0,\dots,T\}$ on $\mathcal{L}_T$. The proof of that lemma
uses improve-or-localize applied to the entire window --- a global descent
identity. Since we are conditioning on $\mathcal{L}_T$ and $k+1 \le T$,
the bound \eqref{eq:confine} applies at index $k+1$ with no further work:
\[
\|u_{k+1}\| \;\le\; S+\Rstar \quad\text{by \eqref{eq:confine} at }j=k+1.
\]
Note that this step does \emph{not} use any inductive hypothesis at $k$.
Confinement is not propagated step-by-step. \Cref{rem:confinement-needs-imploc}
explains why: the projection $|c_{k+1}| \le \|u_{k+1}\|$ only recovers the
$\vmin$-component, so step-local control of $|c_{k+1}|$ alone cannot
bound $\|u_{k+1}\|$.

\subparagraph{Setting up the recurrence.}
The recurrence \eqref{eq:c-lower}, evaluated at index $k$, reads
\begin{equation}\label{eq:c-recurrence-bis}
|c_{k+1}| \;\ge\; (1+\alpha_{k}\epsH)\,|c_k|
   \;-\; \alpha_{k}\,\epsg
   \;-\; \alpha_{k}\,\beta_k\,|\eta_k|
   \;-\; \alpha_{k}\,\|\tilde e_k\|.
\end{equation}
We bound each of the three negative terms by a fraction $\theta_\bullet$ of
the geometric driver $\alpha_{k}\epsH|c_k|$. The key observation is that the
constraints (C-snap), (C-err), (C-mom) are $\alpha_k$-free: each absorption is
proved \emph{before} multiplying by the stepsize, and is then multiplied by
the \emph{same} $\alpha_{k}$ on both sides, so variable stepsizes cost
nothing here.

\subparagraph{Bound on snapshot forcing.}
\emph{Claim.} $\alpha_k\epsg \le \theta''\alpha_k\epsH\,|c_k|$.

\emph{Tools.} (C-snap), inductive hypothesis (iii) at $k$.

\emph{Argument.}
\begin{align*}
\alpha_k\,\epsg
&\;\le\; \alpha_k\,\theta''\,\epsH\,\sdel &&\text{(multiply (C-snap) by }\alpha_k\text{)} \\
&\;\le\; \alpha_k\,\theta''\,\epsH\,|c_k| &&\text{(by (iii): } |c_k| \ge |c_0| \ge \sdel \text{)}.
\end{align*}

\subparagraph{Bound on quadratic-model error.}
\emph{Claim.} $\alpha_k\|\tilde e_k\| \le \theta\alpha_k\epsH\,|c_k|$.

\emph{Tools.} \Cref{lem:graderror} (eq.~\eqref{eq:graderror-snap}),
inductive hypothesis (i) at $k$, (C-err), inductive hypothesis (iii) at $k$.

\emph{Argument.} By \Cref{lem:graderror} applied at $x_k$,
$\|\tilde e_k\| \le \tfrac{\Ltwo}{2}\|u_k\|^2$. The IH (i) at $k$ gives
$\|u_k\| \le S+\Rstar$, so
\[
\|\tilde e_k\| \;\le\; \tfrac{\Ltwo}{2}(S+\Rstar)^2.
\]
Now (C-err) bounds the right-hand side by $\theta\epsH \sdel$, and (iii) at $k$
upgrades $\sdel$ to $|c_k|$:
\[
\|\tilde e_k\| \;\le\; \tfrac{\Ltwo}{2}(S+\Rstar)^2 \;\le\; \theta\,\epsH\,\sdel \;\le\; \theta\,\epsH\,|c_k|.
\]
Multiplying by $\alpha_k$:
\[
\alpha_k\,\|\tilde e_k\| \;\le\; \theta\,\alpha_k\,\epsH\,|c_k|.
\]

\subparagraph{Bound on the momentum term.}
\emph{Claim.} $\alpha_k\,\beta_k|\eta_k| \le \theta'\,\alpha_k\,\epsH\,|c_k|$.

\emph{Tools.} The algorithm's restart test (\Cref{algo:PNCG}, with $q=1$),
gradient upper bound \eqref{eq:barg-def}, gradient floor (ii) at indices
$k-1$ and $k$, (C-mom), IH (iii) at $k$.

\emph{Argument.} First bound $|\eta_k|$:
\[
|\eta_k| \;=\; |\inner{d_{k-1}}{\vmin}| \;\le\; \|d_{k-1}\|.
\]
By the algorithm's restart test (with $q = 1$, $\kappa \ge 1$),
$\|d_{k-1}\| \le \kappa\|g_{k-1}\|$ regardless of whether restart fired:
\begin{itemize}
\item If restart fired at $k-1$: $d_{k-1} = -g_{k-1}$, so $\|d_{k-1}\| = \|g_{k-1}\| \le \kappa\|g_{k-1}\|$ (using $\kappa \ge 1$).
\item If restart did not fire: the contrapositive of the restart condition gives $\|d_{k-1}\| < \kappa\|g_{k-1}\|$.
\end{itemize}
Combined with $\|g_{k-1}\| \le \gup$ from \eqref{eq:barg-def}:
\[
|\eta_k| \;\le\; \kappa\,\|g_{k-1}\| \;\le\; \kappa\,\gup.
\]

For $\beta_k$, the FR formula together with the upper bound
\eqref{eq:barg-def} at index $k$ and the gradient floor at index $k-1$ (which
we have by the strong IH (ii)) gives
$\beta_k = \|g_k\|^2/\|g_{k-1}\|^2 \le \gup^2/\glow^2 = \bbeta$.
Then
\[
\beta_k\,|\eta_k| \;\le\; \bbeta\,\kappa\,\gup.
\]
By (C-mom), this is at most $\theta'\epsH \sdel$; by (iii) at $k$,
$\sdel \le |c_k|$, so
\[
\beta_k\,|\eta_k| \;\le\; \bbeta\,\kappa\,\gup \;\le\; \theta'\,\epsH\,\sdel \;\le\; \theta'\,\epsH\,|c_k|.
\]
Multiplying by $\alpha_k$:
\[
\alpha_k\,\beta_k\,|\eta_k| \;\le\; \theta'\,\alpha_k\,\epsH\,|c_k|.
\]

\emph{Edge case at $k = 0$.} The convention $d_{-1} = 0$ encoding the
post-perturbation momentum reset gives $\eta_0 = \inner{d_{-1}}{\vmin} = 0$,
so the bound holds trivially at $k=0$ without invoking (C-mom) (and without a
floor at index $-1$).

\subparagraph{Growth at $k+1$.}
Substituting the three bounds into \eqref{eq:c-recurrence-bis}:
\begin{align*}
|c_{k+1}|
&\;\ge\; (1+\alpha_{k}\epsH)\,|c_k|
   - \theta''\alpha_{k}\epsH\,|c_k|
   - \theta'\alpha_{k}\epsH\,|c_k|
   - \theta\alpha_{k}\epsH\,|c_k| \\
&\;=\; \bigl(1 + \gamma\,\alpha_{k}\,\epsH\bigr)\,|c_k|,
\end{align*}
using $\gamma = 1 - \theta - \theta' - \theta''$. Combined with IH (iii) at
$k$ and $\alpha_{k}\ge\amin$:
\[
|c_{k+1}|
\;\ge\; \prod_{j=0}^{k}\bigl(1+\gamma\alpha_{j}\epsH\bigr)\,|c_0|
\;\ge\; (1 + \gamma\,\amin\,\epsH)^{k+1}\,|c_0|,
\]
with $|c_0| \ge \sdel$ on $\mathcal{G}_0$ (preserved from the base case).

\subparagraph{Floor at $k+1$.}
\emph{Claim.} $\|g_{k+1}\| \ge \glow$.

\emph{Tools.} \Cref{lem:graderror} at $x_{k+1}$, trigger bound on
$\|g_{\snap}\|$, confinement (i) at $k+1$ (just established), growth (iii) at $k+1$
(just established), definition of $\glow$ \eqref{eq:glow-def}.

\emph{Argument.} \Cref{lem:graderror} gives
$\grad f(x_{k+1}) = g_{\snap} + H\,u_{k+1} + \tilde e_{k+1}$ where
$H = \hess f(\snap)$. Projecting on $\vmin$:
\[
\inner{g_{k+1}}{\vmin}
\;=\; \inner{g_{\snap}}{\vmin}
   + \lambda_{\min}\,c_{k+1}
   + \inner{\tilde e_{k+1}}{\vmin}.
\]
Since $\|g_{k+1}\| \ge |\inner{g_{k+1}}{\vmin}|$ (Cauchy--Schwarz, $\vmin$
unit vector), the reverse triangle inequality gives
\[
\|g_{k+1}\|
\;\ge\; |\lambda_{\min}|\cdot|c_{k+1}|
       - \|g_{\snap}\|
       - \|\tilde e_{k+1}\|.
\]
Now plug in:
\begin{itemize}
\item $|\lambda_{\min}| \ge \epsH$ (trigger condition);
\item $|c_{k+1}| \ge \sdel$ (growth (iii) at $k+1$, combined with $|c_0| \ge \sdel$);
\item $\|g_{\snap}\| \le \epsg$ (trigger condition);
\item $\|\tilde e_{k+1}\| \le \tfrac{\Ltwo}{2}\|u_{k+1}\|^2 \le \tfrac{\Ltwo}{2}(S+\Rstar)^2$ (\Cref{lem:graderror} and confinement (i) at $k+1$).
\end{itemize}
This yields
\[
\|g_{k+1}\|
\;\ge\; \epsH\,\sdel - \epsg - \tfrac{\Ltwo}{2}(S+\Rstar)^2
\;=\; \glow,
\]
by the definition \eqref{eq:glow-def}, which completes the induction.
\end{proof}

\begin{remark}[Why step-local confinement fails]\label{rem:confinement-needs-imploc}
The projection $|c_{k+1}|\le\norm{u_{k+1}}$ recovers only the $\vmin$-component
of the displacement; bounding that component does not control the orthogonal
components. Confinement at step $k+1$ therefore cannot be deduced from the
recurrence~\eqref{eq:c-lower} alone --- it requires the global descent identity
(\Cref{lem:improve}), which enters through the conditioning on $\mathcal{L}_T$
in~\eqref{eq:LT}.
\end{remark}


% =============================================================================
\subsection{Step VIII --- Escape theorem}\label{sub:escape}
% =============================================================================

\paragraph{Rationale.}
Using invariant~(iii) of the induction we show that $|c_{T}|$ exceeds
$(1+\gamma\alpha_k\epsH)^{T}\sdel$; using invariant~(i) we know
$|c_{T}|\le\norm{u_{T}}\le S+\Rstar$. Choosing $T$ so that the first quantity
exceeds the second produces a contradiction, hence
$\mathcal{G}_{0}\cap\mathcal{L}_{T}=\emptyset$: on the good perturbation event,
the localization case is impossible, so the function value must drop by at
least $F$ within the window.

\begin{theorem}[Single-trajectory saddle escape]\label{thm:escape-single}
Define
\begin{equation}\label{eq:Phi-def}
\Phi \;:=\; \log\!\Bigl(\frac{S+R^{*}}{\sdel}\Bigr)
\end{equation}
and choose the window length
\begin{equation}\label{eq:Tchoice}
T \;:=\; \left\lceil
\frac{\Phi}{\log\!\bigl(1+\gamma\,\amin\,\epsH\bigr)}
\right\rceil + 1.
\end{equation}
Under the assumptions of \Cref{lem:induction},
\[
\Pr\Bigl(\,f(x_{0}) - \min_{0\le k\le T}f(x_{k}) \;\ge\; F\,\Bigr)
\;\ge\; \Pr(\mathcal{G}_{0}) \;\ge\; 1 - \delta.
\]
\end{theorem}

\begin{proof}
We show $\mathcal{G}_{0}\cap\mathcal{L}_{T} = \emptyset$. Suppose, for
contradiction, that a sample point lies in
$\mathcal{G}_{0}\cap\mathcal{L}_{T}$. By \Cref{lem:induction}(iii),
\[
|c_{T}| \;\ge\; (1+\gamma\,\amin\,\epsH)^{T}\,\sdel .
\]
Since $\lceil y\rceil+1>y$ for every $y$, the choice~\eqref{eq:Tchoice} gives
\[
T\,\log(1+\gamma\,\amin\,\epsH) \;>\; \Phi,
\qquad\text{hence}\qquad
(1+\gamma\,\amin\,\epsH)^{T}\,\sdel \;>\; e^{\Phi}\sdel \;=\; S+R^{*},
\]
so $|c_{T}|>S+R^{*}$. On the other hand, by \Cref{lem:induction}(i) and
Cauchy--Schwarz, $|c_{T}|\le\norm{u_{T}}\le S+R^{*}$. Contradiction.
Therefore on $\mathcal{G}_{0}$ the event $\mathcal{L}_{T}$ fails, i.e.\ the
function value drops by at least $F$ within the escape window; since
$\Pr(\mathcal{G}_{0})\ge1-\delta$ by \Cref{lem:anticoncentration}, the
conclusion follows.
\end{proof}

\section{Setting and the constraint system}\label{sec:setting}

Throughout, $L_1$ is the gradient Lipschitz constant, $L_2$ the Hessian
Lipschitz constant, $\epsg,\epsH>0$ the tolerances, $d$ the dimension,
$\kappa$ the momentum constant from the restart condition, $\delta\in(0,1)$
the per-window failure probability, and $\LL := f(x_0^{\mathrm{alg}}) -
f_{\mathrm{low}}$ the total decrease budget. Define
\begin{equation}\label{eq:defs}
a \;:=\; \sqrt{\frac{\pi}{2(d+2)}}, \qquad
\Rstar \;:=\; \frac{a\,\epsH\,\delta}{L_2}, \qquad
\sdel \;:=\; \frac{a^2\,\epsH\,\delta^2}{L_2},
\end{equation}
the perturbation radius and the guaranteed initial projection on the escape
direction (probability $\ge 1-\delta$). Direct computation from
\eqref{eq:defs} gives the exact identities
\begin{equation}\label{eq:identities}
\frac{\Rstar}{\sdel} \;=\; \frac{1}{a\,\delta}, \qquad
L_2\,(\Rstar)^2 \;=\; \epsH\,\sdel .
\end{equation}

The escape-window analysis (Sections~7.2--7.5 of the main document) is valid
provided the parameters $S$ (localization radius from the post-perturbation
start), $T$ (window length), $F$ (per-window decrease threshold) and the
splits $\theta,\theta',\theta'',\gamma>0$ satisfy the following system, in
which $\gup := \epsg + L_1(S+\Rstar)$ and
$\glow := \epsH\sdel - \epsg - \tfrac{L_2}{2}(S+\Rstar)^2$ are the upper and
lower gradient envelopes over the window, $\alpha_k$ is the in-window stepsize satisfying $\alpha_k \in [\amin, 1/\Lone]$ with $\amin = \sigma / (\Lone \kappa^2)$, and $\bbeta$ the momentum-coefficient envelope
($\bbeta \le \gup^2/\glow^2$ for Fletcher--Reeves):
\begin{alignat}{2}
&\text{(C-snap)}: \quad && \epsg \;\le\; \theta''\,\epsH\,\sdel,
\label{eq:Csnap}\\
&\text{(C-err)}: \quad && \frac{L_2}{2}\,(S+\Rstar)^2 \;\le\;
\theta\,\epsH\,\sdel, \label{eq:Cerr}\\
&\text{(C-mom)}: \quad && \bbeta\,\kappa\,\gup \;\le\;
\theta'\,\epsH\,\sdel, \label{eq:Cmom}\\
&\text{(C-win)}: \quad && T \;\ge\;
\frac{\log\!\big((S+\Rstar)/\sdel\big)}{\log(1+\gamma\,\amin\,\epsH)},
\label{eq:Cwin}
\end{alignat}
together with the budget constraint
\begin{equation}\label{eq:budget}
\theta + \theta' + \theta'' + \gamma \;\le\; 1 .
\end{equation}

% NOTE: outside the perturbation window, the stepsize remains a free
% placeholder. Inside the window alpha_k in {sigma/(L1 kappa^2), 1/L1}
% (Assumption \ref{ass:pq}); the growth rate is governed by alpha_min, and the
% improve-or-localize constant is 2/L1 uniformly, independently of alpha_k.
% See Remark~\ref{rem:optalpha}.

Define the fixed quantity
\begin{equation}\label{eq:thetamin}
\thetamin \;:=\; \frac{\epsg}{\epsH\,\sdel}
\qquad\text{(so (C-snap) reads } \theta'' \ge \thetamin\text{)},
\end{equation}
and the dimensionless \emph{regime parameter}
\begin{equation}\label{eq:lambda}
\lambda \;:=\; \frac{L_1\,\Rstar}{\epsH\,\sdel}
\;\overset{\eqref{eq:identities}}{=}\; \frac{L_1}{a\,\delta\,\epsH}
\;=\; \frac{L_1}{\delta\,\epsH}\sqrt{\frac{2(d+2)}{\pi}} .
\end{equation}
Small $\lambda$ means the perturbation geometry is large relative to the
gradient-Lipschitz scale; ``deep regime'' below means $\lambda \to 0$.

Finally, normalize the localization radius (we write $\varsigma$, reserving
$\sigma$ for the algorithm's sufficient-descent constant, which enters Step~IX
through $\amin = \sigma/(L_1\kappa^2)$):
\begin{equation}\label{eq:sigma}
\varsigma \;:=\; \frac{L_1\,S}{\epsH\,\sdel}, \qquad\text{so that}\qquad
\frac{L_1\,(S+\Rstar)}{\epsH\,\sdel} \;=\; \varsigma + \lambda .
\end{equation}

\section{The two envelope bounds on $S$}\label{sec:envelopes}

\begin{lemma}[Structural identity and lower bound on $\theta$]
\label{lem:structural}
For every $S \ge 0$, \emph{(C-err)} implies $\theta \ge 1/2$, with strict
inequality whenever $S>0$.
\end{lemma}

\begin{proof}
By \eqref{eq:identities}, $\tfrac{L_2}{2}(\Rstar)^2 =
\tfrac12\,\epsH\sdel$ exactly. Since $S\ge 0$,
$\tfrac{L_2}{2}(S+\Rstar)^2 \ge \tfrac{L_2}{2}(\Rstar)^2 =
\tfrac12\epsH\sdel$, so \eqref{eq:Cerr} forces
$\theta\,\epsH\sdel \ge \tfrac12\epsH\sdel$, i.e.\ $\theta\ge 1/2$.
If $S>0$ the first inequality is strict.
\end{proof}

\begin{lemma}[Error envelope (B-err)]\label{lem:Berr}
\emph{(C-err)} is equivalent to
\begin{equation}\label{eq:Berr}
S + \Rstar \;\le\; \sqrt{2\theta}\,\Rstar,
\qquad\text{i.e.}\qquad
\varsigma \;\le\; \serr(\theta) \;:=\; \big(\sqrt{2\theta}-1\big)\,\lambda .
\end{equation}
\end{lemma}

\begin{proof}
By \eqref{eq:identities},
$\tfrac{L_2}{2}(S+\Rstar)^2 \le \theta\epsH\sdel
\iff (S+\Rstar)^2 \le 2\theta(\Rstar)^2
\iff S+\Rstar \le \sqrt{2\theta}\,\Rstar$ (both sides nonnegative).
Multiplying by $L_1/(\epsH\sdel)$ and using \eqref{eq:lambda},
\eqref{eq:sigma} gives $\varsigma+\lambda \le \sqrt{2\theta}\,\lambda$.
\end{proof}

\begin{lemma}[Momentum envelope (B-mom), Fletcher--Reeves]\label{lem:Bmom}
Suppose \emph{(C-err)} holds and $B := 1-\theta-\thetamin > 0$. Then a
sufficient condition for \emph{(C-mom)} with $\bbeta \le \gup^2/\glow^2$ is
\begin{equation}\label{eq:Bmom}
\varsigma \;\le\; \smom(\theta,\theta')
\;:=\; \Big(\frac{\theta'\,B^2}{\kappa}\Big)^{1/3} - \thetamin - \lambda .
\end{equation}
\end{lemma}

\begin{proof}
Using $\epsg = \thetamin\,\epsH\sdel$ (definition \eqref{eq:thetamin}) and
\eqref{eq:Cerr},
\begin{equation}\label{eq:glow-lb}
\glow \;=\; \epsH\sdel - \epsg - \tfrac{L_2}{2}(S+\Rstar)^2
\;\ge\; \epsH\sdel\,\big(1 - \thetamin - \theta\big)
\;=\; B\,\epsH\,\sdel .
\end{equation}
Since $x \mapsto \kappa\gup^3/x^2$ is decreasing on $x>0$, the lower bound
\eqref{eq:glow-lb} gives the \emph{upper} bound
$\bbeta\kappa\gup \le \kappa\gup^3/\glow^2 \le
\kappa\gup^3/(B\epsH\sdel)^2$ --- note the direction: bounding $\glow$
from \emph{below} is what makes the substitution sufficient rather than
merely necessary. It therefore suffices for \eqref{eq:Cmom} that
\[
\frac{\kappa\,\gup^3}{B^2(\epsH\sdel)^2} \;\le\; \theta'\,\epsH\,\sdel
\quad\Longleftrightarrow\quad
\gup \;\le\; \Big(\frac{\theta' B^2}{\kappa}\Big)^{1/3}\epsH\,\sdel .
\]
With $\gup = \epsg + L_1(S+\Rstar) =
\big(\thetamin + \varsigma + \lambda\big)\epsH\sdel$ by \eqref{eq:thetamin},
\eqref{eq:sigma}, this is exactly \eqref{eq:Bmom}.
\end{proof}

The admissible localization radius is therefore
\begin{equation}\label{eq:min}
\boxed{\;\varsigma \;\le\; \min\big\{\serr(\theta),\;
\smom(\theta,\theta')\big\}.\;}
\end{equation}
% NOTE: the previous resolution treated (C-mom) as the unique binding
% constraint and evaluated S from saturated (B-mom) at theta = 1/2. By
% Lemma~\ref{lem:Berr}, serr(1/2) = 0: that point is infeasible for any S > 0.
% See Remark~\ref{rem:previous}.

\begin{lemma}[Necessary regime condition]\label{lem:regime-nec}
If $\varsigma > 0$ is admissible for some splits satisfying \eqref{eq:budget},
then
\begin{equation}\label{eq:regime-nec}
\lambda + \thetamin \;<\; \frac{B}{\kappa^{1/3}}
\;\le\; \frac{1}{2\,\kappa^{1/3}} .
\end{equation}
\end{lemma}

\begin{proof}
$\smom > 0$ requires $(\theta' B^2/\kappa)^{1/3} > \thetamin+\lambda$.
Since $\theta' \le B$ by \eqref{eq:budget} ($\theta'+\gamma \le B$ and
$\gamma>0$), $(\theta' B^2/\kappa)^{1/3} \le B/\kappa^{1/3}$, and
$B = 1-\theta-\thetamin \le 1/2$ by Lemma~\ref{lem:structural}.
\end{proof}

Unpacking \eqref{eq:lambda}, condition \eqref{eq:regime-nec} requires
$\epsH\,\delta > 2\kappa^{1/3}L_1\sqrt{2(d+2)/\pi}$: the regime restriction
of the single-trajectory analysis. It is \emph{unchanged} by the corrected
optimization below; only the value of $S$ inside the regime changes.

\section{Optimization of the splits}\label{sec:optimization}

\subsection{The objective and its exact surrogate}

The quantities derived in Section~\ref{sec:derived} give
$\Niter \le 2T^2\LL/(L_1 S^2) + T$ with
$T \le \Phi/(\gamma\,\amin\,\epsH) + \Phi + 2$ and
$\Phi = \log((S+\Rstar)/\sdel)$. We therefore \emph{maximize}
\begin{equation}\label{eq:objective}
J \;:=\; \gamma\,\varsigma
\end{equation}
over the splits, subject to \eqref{eq:min}, \eqref{eq:budget},
$\theta\ge 1/2$, $\theta'' \ge \thetamin$.

\begin{remark}[Exactness of the surrogate]\label{rem:surrogate}
$\Niter$ and $J$ are related by
$\Niter \lesssim T^2/S^2 \propto 1/(\gamma\varsigma)^2 = 1/J^2$ up to three
explicitly bounded discrepancies: (i) $\log(1+x)\in[x/(1+x),x]$, so under
$\gamma\,\amin\,\epsH\le 1$ the bound on $T$ is exact within a factor $2$,
hence $\Niter$ within a factor $4$; (ii) $\Phi$ depends on the splits only
through $\log(S+\Rstar)$, a logarithmic variation; (iii) the additive
terms $+2$, $+T$. Consequently, maximizing $J$ minimizes the
$\Niter$ bound up to an explicit absolute factor; combined with
Proposition~\ref{prop:nearopt} this yields
$\Niter^{\mathrm{canonical}} \le 16\,\Niter^{\mathrm{opt}}$
under \eqref{eq:Rprime}. No claim sharper than this factor is made.
\end{remark}

\subsection{Structure of the optimum}

\begin{lemma}[$\theta''$ is pinned]\label{lem:thetapp}
At any maximizer of $J$ with $J>0$, $\theta'' = \thetamin$.
\end{lemma}

\begin{proof}
$\theta''$ enters $J$ only through the budget \eqref{eq:budget}: neither
$\serr$ nor $\smom$ depends on $\theta''$ (the quantity $\thetamin$ in
\eqref{eq:Bmom} is the fixed problem datum \eqref{eq:thetamin}, not the split
$\theta''$, and $B$ is defined with $\thetamin$). If $\theta'' > \thetamin$,
replace $(\theta'',\gamma)$ by $(\thetamin,\ \gamma + \theta''-\thetamin)$:
all constraints still hold and $J$ strictly increases.
\end{proof}

In view of Lemma~\ref{lem:thetapp} we set $\theta'' = \thetamin$, write
\begin{equation}\label{eq:MB}
M \;:=\; 1 - \thetamin, \qquad B \;=\; M - \theta,
\qquad \gamma \;=\; B - \theta' ,
\end{equation}
and saturate the budget \eqref{eq:budget} (any slack is given to $\gamma$,
which strictly increases $J$). The remaining variables are
$(\theta,\theta') \in [1/2, M)\times(0, B)$.

\begin{lemma}[Balance at the optimum]\label{lem:balance}
Assume $J^\star := \sup J > 0$. Then the supremum is attained, and at every
maximizer
\begin{equation}\label{eq:balance}
\serr(\theta) \;=\; \smom(\theta,\theta'),
\qquad\text{i.e.}\qquad
\Big(\frac{\theta' B^2}{\kappa}\Big)^{1/3}
\;=\; \thetamin + \sqrt{2\theta}\,\lambda .
\end{equation}
\end{lemma}

\begin{proof}
$J$ is continuous on the compact set
$\{\theta\in[1/2,M],\ \theta'\in[0,M-\theta]\}$ (with
$\varsigma := \min\{\serr,\smom\}^+$), so a maximizer exists.

\emph{Case $\smom > \serr$ at a maximizer.} Then
$\smom > \serr \ge 0$ forces $\theta' > 0$. Decrease $\theta'$ by
$\eta>0$ and give $\eta$ to $\gamma$. By continuity of $\smom$ in
$\theta'$, for $\eta$ small enough still $\smom \ge \serr$, so
$\varsigma = \serr$ is unchanged while $\gamma$ strictly increases:
$J$ strictly increases, contradiction.

\emph{Case $\smom < \serr$ at a maximizer.} Then $J>0$ forces
$\smom>0$, and $\serr > \smom > 0$ forces $\sqrt{2\theta}>1$, i.e.\
$\theta > 1/2$ strictly. Decrease $\theta$ by $\eta>0$ and give $\eta$ to
$\gamma$ (keeping $\theta'$): $B$ increases, so $\smom$ strictly increases;
$\serr$ decreases continuously, so for $\eta$ small enough still
$\serr > \smom$. Hence $\varsigma = \smom$ strictly increased and $\gamma$
strictly increased: $J$ strictly increases, contradiction.

Therefore $\serr = \smom$ at every maximizer; the displayed equivalent form
follows by adding $\thetamin+\lambda$ to $\serr = (\sqrt{2\theta}-1)\lambda$
and cubing.
\end{proof}

\subsection{Exact upper bound and the canonical splits}

\begin{lemma}[Exact upper bound on the objective]\label{lem:UB}
Let $M > 1/2$ and
\begin{equation}\label{eq:tstar}
\tstar \;:=\; \frac{1+\sqrt{1+6M}}{3} \;\in\; \big(1,\sqrt{2M}\,\big).
\end{equation}
Then for \emph{every} feasible choice of splits,
\begin{equation}\label{eq:UB}
J \;\le\; \lambda\,\tstar\,(\tstar-1)^2 .
\end{equation}
\end{lemma}

\begin{proof}
For any feasible point, $J = \gamma\sigma \le \gamma\,\serr \le
B\,\serr = (M-\theta)\big(\sqrt{2\theta}-1\big)\lambda$, using
$\gamma = B - \theta' \le B$. Substitute $t := \sqrt{2\theta} \in
[1,\sqrt{2M}]$ and define
\[
h(t) \;:=\; \Big(M - \tfrac{t^2}{2}\Big)(t-1),
\qquad
h'(t) \;=\; -\tfrac{3}{2}t^2 + t + M,
\qquad
h''(t) \;=\; -3t+1 < 0 \ \ (t\ge 1).
\]
$h$ is strictly concave on $[1,\sqrt{2M}]$ with
$h'(1) = M - \tfrac12 > 0$ and
$h'(\sqrt{2M}) = \sqrt{2M} - 2M < 0$ (since $M>1/2$), so $h$ has a unique
interior maximizer, the positive root of $h'$, which is \eqref{eq:tstar}.
At $\tstar$ the first-order condition gives the exact identity
\begin{equation}\label{eq:FOC}
M \;=\; \frac{3\tstar^2 - 2\tstar}{2}
\qquad\Longrightarrow\qquad
M - \frac{\tstar^2}{2} \;=\; \tstar^2 - \tstar \;=\; \tstar(\tstar-1),
\end{equation}
hence $h(\tstar) = \tstar(\tstar-1)^2$, which is \eqref{eq:UB}.
\end{proof}

\begin{definition}[Canonical splits]\label{def:canonical}
With $\tstar$ as in \eqref{eq:tstar}, set
\begin{equation}\label{eq:canonical}
\theta_\star := \frac{\tstar^2}{2}, \qquad
\theta''_\star := \thetamin, \qquad
B_\star := M - \theta_\star
\overset{\eqref{eq:FOC}}{=} \tstar(\tstar-1), \qquad
\theta'_\star := \frac{\kappa\,\big(\thetamin +
\tstar\lambda\big)^3}{B_\star^2}, \qquad
\gamma_\star := B_\star - \theta'_\star ,
\end{equation}
and take $\sigma$ at its admissible maximum,
$\sigma_\star := (\tstar - 1)\,\lambda$, i.e.
\begin{equation}\label{eq:Sstar}
S \;=\; (\tstar-1)\,\Rstar
\;=\; (\tstar - 1)\,\sqrt{\frac{\pi}{2(d+2)}}\;\frac{\epsH\,\delta}{L_2},
\qquad
S + \Rstar \;=\; \tstar\,\Rstar .
\end{equation}
\end{definition}

\begin{lemma}[Feasibility of the canonical splits]\label{lem:feasible}
Assume $\thetamin < 1/2$ (tolerance compatibility, equivalently
$\epsg < \tfrac12\,\epsH\sdel$) and the \emph{regime restriction}
\begin{equation}\label{eq:Rprime}
\thetamin + \tstar\,\lambda \;\le\; \frac{B_\star}{(2\kappa)^{1/3}} .
\tag{R$'$}
\end{equation}
Then the canonical splits are feasible: \emph{(C-snap)}, \emph{(C-err)},
\emph{(C-mom)} all hold with $S$ as in \eqref{eq:Sstar},
$\theta'_\star \le B_\star/2$, and $\gamma_\star \ge B_\star/2 > 0$.
Moreover \emph{(C-err)} and \emph{(C-mom)} are \emph{saturated} (hold with
equality in the envelope bounds \eqref{eq:Berr}, \eqref{eq:Bmom}).
\end{lemma}

\begin{proof}
$\thetamin<1/2$ gives $M>1/2$, so $\tstar\in(1,\sqrt{2M})$ and
$\theta_\star = \tstar^2/2 \in (1/2, M)$, $B_\star>0$.
\eqref{eq:Rprime} cubed reads
$\kappa(\thetamin+\tstar\lambda)^3 \le B_\star^3/2$, i.e.\
$\theta'_\star \le B_\star/2$, whence $\gamma_\star \ge B_\star/2 > 0$ and
\eqref{eq:budget} holds with equality by construction.
(C-err): by \eqref{eq:Sstar} and \eqref{eq:identities},
$\tfrac{L_2}{2}(S+\Rstar)^2 = \tfrac{\tstar^2}{2}\,\epsH\sdel =
\theta_\star\,\epsH\sdel$ --- equality.
(C-mom): $\gup = (\thetamin + \tstar\lambda)\,\epsH\sdel$ by
\eqref{eq:sigma}, and by \eqref{eq:glow-lb},
$\glow \ge B_\star\,\epsH\sdel$, so
\[
\bbeta\,\kappa\,\gup \;\le\; \frac{\kappa\,\gup^3}{\glow^2}
\;\le\; \frac{\kappa\,(\thetamin+\tstar\lambda)^3}{B_\star^2}\;\epsH\sdel
\;=\; \theta'_\star\,\epsH\,\sdel ,
\]
with equality in the last step by definition of $\theta'_\star$:
\eqref{eq:Cmom} holds and the envelope \eqref{eq:Bmom} is saturated, i.e.\
$\smom(\theta_\star,\theta'_\star) = \sigma_\star = \serr(\theta_\star)$,
which is the balance equation \eqref{eq:balance}.
\end{proof}

\begin{proposition}[Near-optimality, exact form]\label{prop:nearopt}
Let
\begin{equation}\label{eq:rho}
\rho_\star \;:=\; \frac{\theta'_\star}{B_\star}
\;=\; \kappa\,\Big(\frac{\thetamin + \tstar\lambda}{\tstar(\tstar-1)}
\Big)^{3} .
\end{equation}
Under the assumptions of Lemma~\ref{lem:feasible},
\begin{equation}\label{eq:nearopt}
J_{\mathrm{can}} \;=\; \gamma_\star\,\sigma_\star
\;=\; (1-\rho_\star)\,\lambda\,\tstar(\tstar-1)^2
\;\ge\; (1-\rho_\star)\;\sup J ,
\qquad \rho_\star \le \tfrac12 .
\end{equation}
In particular $J_{\mathrm{can}} \ge \tfrac12 \sup J$ always under
\eqref{eq:Rprime}, and $\rho_\star \to 0$ as
$\lambda,\thetamin \to 0$: deep in the regime the canonical splits are
asymptotically exactly optimal.
\end{proposition}

\begin{proof}
$J_{\mathrm{can}} = (B_\star - \theta'_\star)(\tstar-1)\lambda =
(1-\rho_\star)B_\star(\tstar-1)\lambda =
(1-\rho_\star)\,\tstar(\tstar-1)^2\lambda$ by \eqref{eq:FOC}. Combine with
the upper bound \eqref{eq:UB}, valid for every feasible point. The bound
$\rho_\star\le 1/2$ is \eqref{eq:Rprime} cubed.
\end{proof}

\begin{remark}[Why the previous resolution is invalid]\label{rem:previous}
The earlier derivation fixed $\theta = 1/2$ (``with infinitesimal margin''),
declared (C-mom) the binding constraint, and evaluated $S$ from saturated
(B-mom), obtaining $S + \Rstar = B\,\epsH\sdel/\big(L_1(4\kappa)^{1/3}\big)$
and hence $S \propto \delta^2$. By Lemma~\ref{lem:Berr}, however,
$\serr(1/2) = 0$: at $\theta = 1/2$, (C-err) forces $S = 0$, and near
$\theta=1/2$, $S \le (\sqrt{2\theta}-1)\Rstar \approx (\theta-\tfrac12)\Rstar$
scales \emph{linearly} in the margin. The value of $S$ used in the previous
Steps 6--10 violates (C-err) whenever it is positive, deep in the regime
$\lambda \ll 1$ --- which is precisely where it was invoked. The binding
structure is in fact: (C-mom) binds only near the boundary of
\eqref{eq:Rprime}; deep in the regime both constraints bind simultaneously at
the optimum (Lemma~\ref{lem:balance}), with $S = (\tstar-1)\Rstar \propto
\delta$, not $\delta^2$. Consequently the previously stated per-window count
$\Niter \propto L_1^3 L_2^2\,d^2\,\LL/(\epsH^6\delta^4)$ does not correspond
to any feasible instantiation of the constraint system and must be
retracted; the feasible count is Theorem~\ref{thm:main} below, which exceeds
it by the factor $\lambda^{-2}$ (up to absolute constants) inside the
regime, the two coinciding at the boundary of \eqref{eq:Rprime}.
\end{remark}

\begin{remark}[Anchor point and the radius $S+2\Rstar$]\label{rem:anchor}
% NOTE: action item --- audit the ~24 sites where the confinement radius was
% changed to S + 2R^* and classify each as (a) geometric localization
% relative to the underlying saddle (harmless) or (b) input to (C-err) or to
% the definition of \gup (must be snap-anchored).
All four constraints above are anchored at the \emph{snap point}, where the
trigger certifies $\lambda_{\min}(\nabla^2 f(\mathrm{snap})) \le -\epsH$;
the confinement radius from the snap point is $S + \Rstar$ (perturbation
$\le \Rstar$, in-window travel $\le S$), and this is what enters (C-err) and
$\gup$. If instead (C-err) is stated with radius $S + \nu\Rstar$ for
$\nu \ge 1$ (e.g.\ $\nu = 2$ when anchoring at the underlying saddle, at
distance up to $\Rstar$ from the snap point), the identity of
Lemma~\ref{lem:structural} becomes
$\tfrac{L_2}{2}(\nu\Rstar)^2/(\epsH\sdel) = \nu^2/2$, forcing
$\theta \ge \nu^2/2$; for $\nu = 2$ this gives $\theta \ge 2$, incompatible
with \eqref{eq:budget}: the system is infeasible. The infeasibility cannot
be repaired by shrinking the perturbation radius alone: if the perturbation
radius is $\Rstar/\varrho$ ($\varrho \ge 1$) while the snap-to-saddle
distance remains $\le \Rstar$, then $\sdel$ scales by $1/\varrho$ and the
saddle-anchored identity becomes
$\varrho\,(1+1/\varrho)^2/2 = (\varrho + 2 + 1/\varrho)/2 \ge 2$ for all
$\varrho \ge 1$. Hence saddle-anchoring of (C-err) is structurally
infeasible under the present definitions \eqref{eq:defs}; the analysis must
remain snap-anchored, which is also what the trigger certifies.
\end{remark}

\section{Derived quantities: $S$, $T$, $F$, $\Niter$}\label{sec:derived}

Throughout this section the canonical splits \eqref{eq:canonical} are in
force, $S$ is given by \eqref{eq:Sstar}, and the in-window stepsizes satisfy
$\alpha_k \in \{\sigma/(L_1\kappa^2),\, 1/L_1\} \subset [\amin,\, 1/L_1]$
(\Cref{ass:pq}), with $\amin = \sigma/(L_1\kappa^2)$.

\paragraph{Window length.} Set
\begin{equation}\label{eq:Phi}
\Phi \;:=\; \log\!\Big(\frac{S+\Rstar}{\sdel}\Big)
\;\overset{\eqref{eq:Sstar},\eqref{eq:identities}}{=}\;
\log\!\Big(\frac{\tstar}{a\,\delta}\Big)
\;=\; \log\!\Big(\frac{\tstar}{\delta}\sqrt{\frac{2(d+2)}{\pi}}\Big),
\qquad
T \;:=\; \Big\lceil \frac{\Phi}{\log(1+\gamma_\star\,\amin\,\epsH)}
\Big\rceil + 1,
\end{equation}
so that (C-win) holds by construction; this is the choice
\eqref{eq:Tchoice} of \Cref{thm:escape-single} with $\gamma = \gamma_\star$
and $S+\Rstar = \tstar\Rstar$. Using $\log(1+x)\ge x/(1+x)$ for $x \ge 0$,
\begin{equation}\label{eq:T}
T \;\le\; \frac{\Phi}{\log(1+\gamma_\star\,\amin\,\epsH)} + 2
\;\le\; \frac{(1+\gamma_\star\,\amin\,\epsH)\,\Phi}{\gamma_\star\,\amin\,\epsH} + 2
\;=\; \frac{L_1\,\kappa^2\,\Phi}{\sigma\,\gamma_\star\,\epsH} + \Phi + 2 .
\tag{T}
\end{equation}
% NOTE (flagged corrections to (T)): (1) the denominator is
% log(1+gamma*alpha_min*epsH), not gamma*alpha*epsH (log(1+x) <= x, so the old
% T was too small); (2) the growth rate is only guaranteed at the smallest
% in-window stepsize alpha_min; (3) the old second inequality
% "(L1+gamma*epsH)Phi/(gamma*epsH)+1" was valid only at alpha = 1/L1 and is
% superseded.

\paragraph{Per-window decrease.} By \Cref{lem:localization}, the
localization radius on the no-decrease event is $S = \sqrt{2FT/L_1}$: the
constant $2/L_1$ is the \emph{uniform} improve-or-localize constant of
\Cref{lem:improve}, valid for every in-window step under \Cref{ass:pq}
(restarted steps and non-restarted steps with $\alpha^{\ast}_{\setN}$ are
the only two iteration types). Equivalently, choosing
\begin{equation}\label{eq:F}
F \;:=\; \frac{L_1\,S^2}{2\,T}
\;=\; \frac{(\tstar-1)^2\,\pi\,L_1\,\epsH^2\,\delta^2}
{4\,(d+2)\,T\,L_2^2}
\tag{F}
\end{equation}
makes the localization radius exactly $S$.
% NOTE (flagged correction): the previous version set F = S^2/(2 alpha T),
% inherited from the GD-style bound sum ||dx||^2 <= 2 alpha (f(x_0)-f(x_T)).
% The document's own improve-or-localize lemma gives the constant 2/L1 for
% both iteration types, independently of alpha_k; (F) is now consistent with
% the definition of S in Lemma \ref{lem:localization}.

\paragraph{Iteration count (window iterations).} Lifetime accounting:
the number of perturbation windows satisfies
$\Ntrig \le \LL/F + 1$, hence
\begin{equation}\label{eq:N}
\Niter \;\le\; T\,\Ntrig
\;\le\; \frac{2\,T^2\,\LL}{L_1\,S^2} + T
\;=\; \frac{4\,(d+2)\,T^2\,L_2^2\,\LL}
{\pi\,L_1\,(\tstar-1)^2\,\epsH^2\,\delta^2} + T .
\tag{N}
\end{equation}

\begin{remark}[Role of the in-window stepsizes]\label{rem:optalpha}
By \eqref{eq:T}, $T \propto 1/\amin$ to leading order, so \eqref{eq:N}
scales as $\Niter \propto T^2 \propto \amin^{-2} = (L_1\kappa^2/\sigma)^2$:
the price of taking genuine NCG steps inside the window, rather than
gradient steps at $1/L_1$, is the absolute factor $\kappa^4/\sigma^2$.
Since $(\sigma,\kappa)$ are algorithm constants, this does not affect the
dependence on $(\epsH, d, \delta)$. Outside the perturbation window the
stepsize remains a free placeholder, unconstrained by this analysis.
\end{remark}

\begin{corollary}[Explicit count]\label{cor:explicit}
Assume additionally $\gamma_\star\,\epsH \le L_1$ (which implies
$\gamma_\star\,\amin\,\epsH \le \sigma/\kappa^2 \le 1$) and
$a\,\delta \le \tstar/e$ (so $\Phi \ge 1$). Then
\[
T \;\le\; \frac{4\,L_1\,\kappa^2\,\Phi}{\sigma\,\gamma_\star\,\epsH}
\]
and
\begin{equation}\label{eq:Nexplicit}
\Niter \;\le\;
\frac{64\,(d+2)\,\kappa^4\,L_1\,L_2^2\,\Phi^2\,\LL}
{\pi\,\sigma^2\,\gamma_\star^2\,(\tstar-1)^2\,\epsH^4\,\delta^2}
\;+\; \frac{4\,L_1\,\kappa^2\,\Phi}{\sigma\,\gamma_\star\,\epsH} .
\end{equation}
\end{corollary}

\begin{proof}
By \eqref{eq:T}, $T \le L_1\kappa^2\Phi/(\sigma\gamma_\star\epsH) + \Phi + 2$.
Since $\Phi \ge 1$, $\Phi + 2 \le 3\Phi$; since
$\gamma_\star\amin\epsH \le 1$,
$\Phi \le \Phi/(\gamma_\star\amin\epsH)
= L_1\kappa^2\Phi/(\sigma\gamma_\star\epsH)$. Hence
$T \le 4L_1\kappa^2\Phi/(\sigma\gamma_\star\epsH)$. Substituting
$T^2 \le 16L_1^2\kappa^4\Phi^2/(\sigma^2\gamma_\star^2\epsH^2)$ and
$S^2 = (\tstar-1)^2\pi\epsH^2\delta^2/\bigl(2(d+2)L_2^2\bigr)$ into
\eqref{eq:N} gives \eqref{eq:Nexplicit}.
\end{proof}

\section{Union bound over windows}\label{sec:union}

Each window fails (the perturbation misses the escape direction) with
probability at most $\delta$. On the event that all of the first
$\lceil \LL/F\rceil$ windows succeed, the algorithm performs at most
$\Ntrig \le \LL/F + 1$ windows; by the union bound the failure probability
is at most $(\LL/F + 1)\,\delta$. Writing $F = c'\,\delta^2$ with
\begin{equation}\label{eq:cprime}
c' \;:=\; \frac{\pi\,(\tstar-1)^2\,L_1\,\epsH^2}{4\,(d+2)\,T\,L_2^2}
\qquad (\text{which depends on } \delta \text{ only through }
\Phi(\delta) \text{ inside } T),
\end{equation}
the requirement of total failure probability at most
$\Delta \in (0,1)$ reads
\begin{equation}\label{eq:union}
\frac{\LL}{c'(\delta)\,\delta} + \delta \;\le\; \Delta .
\tag{U}
\end{equation}

\begin{lemma}[Admissible interval for $\delta$]\label{lem:admissible}
A sufficient condition for \eqref{eq:union} is
\begin{equation}\label{eq:window}
\max\Big\{\;\delta_{\mathrm{reg}},\;\;
\delta_{\mathrm{u}}\;\Big\}
\;\le\; \delta \;\le\; \frac{\Delta}{2},
\qquad\text{where}\qquad
\delta_{\mathrm{u}} \text{ solves } \delta_{\mathrm{u}} =
\frac{2\,\LL}{c'(\delta_{\mathrm{u}})\,\Delta},
\end{equation}
and $\delta_{\mathrm{reg}}$ is the smallest $\delta$ satisfying the regime
restriction \eqref{eq:Rprime}, namely
\begin{equation}\label{eq:deltareg}
\delta_{\mathrm{reg}} \;=\;
\frac{\tstar\,(2\kappa)^{1/3}\,L_1}
{a\,\epsH\,\big(B_\star - (2\kappa)^{1/3}\,\thetamin\big)},
\qquad\text{requiring}\qquad
B_\star > (2\kappa)^{1/3}\,\thetamin .
\end{equation}
The threshold $\delta_{\mathrm{u}}$ exists and is unique: the map
$\delta \mapsto 2\LL/(c'(\delta)\Delta)$ is strictly decreasing (as
$\delta$ increases, $\Phi$ decreases, $T$ decreases, $c'$ increases), so it
crosses the identity exactly once on $(0,1)$.
The interval \eqref{eq:window} is nonempty if and only if
$\delta_{\mathrm{reg}} \le \Delta/2$ and
$\delta_{\mathrm{u}} \le \Delta/2$; the latter is equivalent to
$c'(\Delta/2)\,\Delta^2 \ge 4\,\LL$.
\end{lemma}

\begin{proof}
If $\delta \le \Delta/2$ and $\delta \ge 2\LL/(c'(\delta)\Delta)$, then
$\LL/(c'\delta) \le \Delta/2$ and \eqref{eq:union} follows by summing.
Monotonicity of $c'$ in $\delta$: $\Phi(\delta) = \log(\tstar/(a\delta))$
is strictly decreasing; $T$ in \eqref{eq:Phi} is nondecreasing in $\Phi$;
$c' \propto 1/T$. The crossing claim follows since the identity map is
strictly increasing from $0$ and the decreasing map is positive.
Nonemptiness: $\delta \ge \delta_{\mathrm{u}}$ for $\delta \le \Delta/2$
is equivalent to $\Delta/2 \ge \delta_{\mathrm{u}}$ by monotonicity, which
unfolds to $c'(\Delta/2)\,(\Delta/2) \ge 2\LL/\Delta$.
\end{proof}

\begin{remark}[Choice of $\delta$ within the admissible interval]
\label{rem:deltachoice}
The bound \eqref{eq:N} is strictly decreasing in $\delta$ (as
$1/\delta^2$ up to the logarithm in $T$), so within \eqref{eq:window} the
iteration bound is minimized at the upper endpoint $\delta = \Delta/2$.
% TODO (interacts with the open gap): the perturbation's function-value jump
% f(x_0) - f(snap) is not yet rigorously bounded in the main document; it
% grows with the perturbation radius R^* \propto \delta. Pinning
% \delta = \Delta/2 maximizes this exposure; until the jump is bounded and
% absorbed into the budget accounting, the conservative choice is the lower
% endpoint of (\ref{eq:window}). The theorem below is stated for an
% arbitrary admissible \delta for this reason.
\end{remark}

\begin{remark}[Structural change versus the previous coupling equation]
\label{rem:coupling}
Because $F \propto \delta^2$ (not $\delta^4$), the self-consistency
relation for $\delta$ is the degree-one crossing of
Lemma~\ref{lem:admissible}, not the previous cubic
$\delta^3 = 2T(4\kappa)^{2/3}L_1L_2^2\LL/(B^2a^4\epsH^4\Delta)$. That
equation inherited the exponent $4$ from the infeasible
$S \propto \delta^2$ and is superseded.
\end{remark}

\section{Main theorem}\label{sec:main}

\begin{theorem}[PR--NCG escape: corrected exact constants]\label{thm:main}
Fix tolerances $\epsg, \epsH > 0$ and a global failure probability
$\Delta \in (0,1)$. Define $a$, $\Rstar$, $\sdel$, $\thetamin$, $\lambda$ by
\eqref{eq:defs}, \eqref{eq:thetamin}, \eqref{eq:lambda};
$M = 1-\thetamin$ and $\tstar$ by \eqref{eq:tstar}; the canonical splits by
\eqref{eq:canonical}; and let $\delta$ be any value in the admissible
interval \eqref{eq:window}. Assume:
\begin{enumerate}
\item \textbf{Tolerance compatibility:} $\thetamin < 1/2$ and
$B_\star > (2\kappa)^{1/3}\thetamin$;
\item \textbf{Regime restriction \eqref{eq:Rprime}:}
$\thetamin + \tstar\lambda \le B_\star/(2\kappa)^{1/3}$
(equivalently $\delta \ge \delta_{\mathrm{reg}}$ of
\eqref{eq:deltareg});
\item \textbf{Union-bound compatibility:} the interval \eqref{eq:window}
is nonempty;
\item the escape-window analysis (\Cref{lem:induction},
\Cref{thm:escape-single}) holds under
\eqref{eq:Csnap}--\eqref{eq:Cwin}, with $p=q=1$ (\Cref{ass:pq}) and
in-window stepsizes $\alpha_k \in \{\sigma/(L_1\kappa^2),\, 1/L_1\}$.
\end{enumerate}
Set $S = (\tstar-1)\Rstar$, $T$ as in \eqref{eq:Phi}, and
$F = L_1 S^2/(2T)$. Then with probability at least $1-\Delta$, every
perturbation window either decreases $f$ by at least $F$ or escapes the
$\epsH$-saddle, the total number of windows is at most $\LL/F + 1$, and the
total number of window iterations satisfies the exact bound \eqref{eq:N};
under the additional assumptions of \Cref{cor:explicit}
($\gamma_\star\epsH \le L_1$, $a\delta \le \tstar/e$),
\[
\Niter \;\le\;
\frac{64\,(d+2)\,\kappa^4\,L_1\,L_2^2\,\Phi^2\,\LL}
{\pi\,\sigma^2\,\gamma_\star^2\,(\tstar-1)^2\,\epsH^4\,\delta^2}
\;+\; \frac{4\,L_1\,\kappa^2\,\Phi}{\sigma\,\gamma_\star\,\epsH},
\qquad
\Phi = \log\!\Big(\frac{\tstar}{\delta}\sqrt{\tfrac{2(d+2)}{\pi}}\Big).
\]
Moreover the canonical splits are optimal for the surrogate objective
$J = \gamma\varsigma$ within the exact factor $(1-\rho_\star) \ge 1/2$ of
Proposition~\ref{prop:nearopt}, hence optimal for the iteration bound
within an explicit absolute factor (Remark~\ref{rem:surrogate}); and by
Lemma~\ref{lem:regime-nec} no choice of splits removes the regime
restriction $\lambda + \thetamin < 1/(2\kappa^{1/3})$, i.e.\
$\epsH\,\delta > 2\kappa^{1/3}L_1\sqrt{2(d+2)/\pi}$.
\end{theorem}

\begin{proof}
Feasibility of the splits and saturation of (C-err), (C-mom) is
Lemma~\ref{lem:feasible}; (C-win) holds by the definition of $T$ in
\eqref{eq:Phi}; (C-snap) holds by \eqref{eq:thetamin} and
$\theta''_\star = \thetamin$. Assumption~4 then yields the per-window
dichotomy with failure probability at most $\delta$ per window
(\Cref{thm:escape-single}). The window count and \eqref{eq:N} are the
lifetime accounting of Section~\ref{sec:derived}; the probability statement
is Lemma~\ref{lem:admissible}. Near-optimality is
Proposition~\ref{prop:nearopt} with Remark~\ref{rem:surrogate}, and the
irremovability of the regime restriction is Lemma~\ref{lem:regime-nec},
which used only \eqref{eq:budget} and Lemma~\ref{lem:structural}, i.e.\
holds for \emph{every} feasible choice of splits.
\end{proof}

\begin{remark}[Scaling summary and comparison]\label{rem:scaling}
Per window (fixed $\delta$), the corrected count scales as
$\Niter = \widetilde{O}\!\big(\kappa^4\sigma^{-2}\,L_1 L_2^2\,(d+2)\,\LL\,
\epsH^{-4}\,\delta^{-2}\big)$, valid only under \eqref{eq:Rprime}; the
factor $\kappa^4/\sigma^2$ is an absolute constant of the algorithm
(\Cref{rem:optalpha}). Against Jin--Netrapalli--Jordan ($\epsH^{-4}$,
polylog $d$, no regime restriction), the remaining gaps are the
$(d+2)\,\delta^{-2}$ factor and the regime restriction itself; both
originate in the single-trajectory anti-concentration constant
$a \propto d^{-1/2}$ and the Fletcher--Reeves envelope $\gup/\glow$ (the
cubic in Lemma~\ref{lem:Bmom}), respectively, and neither is touched by any
choice of splits. Removing them is the object of the coupling extension.
\end{remark}




% \subsection{Step 3c: coupling two trajectories}\label{sub:step3c}
% Following \cite{pmlr-v70-jin17a}, we track two coupled copies of the algorithm and show
% their separation grows. Let $v_{\min}$ be a unit minimal eigenvector of
% $\hess f(\snap)$ and let $\rho > 0$ be a small shift. Start the two copies from
% \begin{equation}\label{eq:couple-init}
% x_{0} = \snap + \xi,
% \qquad
% x_{0}' = \snap + \xi + \rho\, v_{\min},
% \end{equation}
% and define the stacked difference vector
% \begin{equation}\label{eq:wdef}
% w_{k} := \begin{pmatrix} x_{k} - x_{k}' \\ d_{k-1} - d_{k-1}' \end{pmatrix}
%        =: \begin{pmatrix} \Delta x_{k} \\ \Delta d_{k-1} \end{pmatrix},
% \end{equation}
% where $\Delta(\cdot)$ denotes the difference of a quantity between the two copies.

% \paragraph{Difference recursion.}
% From $x_{k+1} = x_{k} + \alpha_{k} d_{k}$ and $x_{k+1}' = x_{k}' + \alpha_{k}' d_{k}'$,
% splitting $\alpha_{k} d_{k}-\alpha_{k}' d_{k}'
% = \alpha_{k}(d_{k}-d_{k}') + (\alpha_{k}-\alpha_{k}')d_{k}'$,
% \begin{equation}\label{eq:dx-rec}
% \Delta x_{k+1} = \Delta x_{k} + \alpha_{k}\,\Delta d_{k} + \Delta\alpha_{k}\, d_{k}'.
% \end{equation}
% Likewise, from $d_{k} = -g_{k} + \beta_{k} d_{k-1}$,
% \begin{equation}\label{eq:dd-rec}
% \Delta d_{k} = -\Delta g_{k} + \beta_{k}\,\Delta d_{k-1} + \Delta\beta_{k}\, d_{k-1}'.
% \end{equation}
% Linearizing the gradient with \Cref{lem:graderror}, $g_{k} = H(x_{k}-\xstar)+e_{k}$ and
% $g_{k}' = H(x_{k}'-\xstar)+e_{k}'$,
% \begin{equation}\label{eq:dg}
% \Delta g_{k} = H\,\Delta x_{k} + (e_{k} - e_{k}').
% \end{equation}

% \paragraph{The error difference is Lipschitz.}
% Let $E(x) := \grad f(x) - H(x-\xstar)$, so $e_{k} = E(x_{k})$, $e_{k}' = E(x_{k}')$. Its
% Jacobian is $J_{E}(x) = \hess f(x) - H = \hess f(x) - \hess f(\xstar)$. On the ball
% $\norm{x-\xstar}\leq S$, \Cref{ass:hesslip} gives
% $\norm{J_{E}(x)} \leq \Ltwo\norm{x-\xstar}\leq \Ltwo S$, so $E$ is $\Ltwo S$-Lipschitz
% there. Parametrizing the segment $y_{s} = x_{k}' + s(x_{k}-x_{k}')$, $s\in[0,1]$ (which
% stays in the ball, assumed convex), the FTC yields
% $E(x_{k})-E(x_{k}') = \int_{0}^{1} J_{E}(y_{s})(x_{k}-x_{k}')\,ds$, hence
% \begin{equation}\label{eq:ediff}
% \norm{e_{k} - e_{k}'} \;\leq\; \int_{0}^{1}\norm{J_{E}(y_{s})}\,\norm{\Delta x_{k}}\,ds
%   \;\leq\; \Ltwo S\,\norm{\Delta x_{k}}.
% \end{equation}

% \paragraph{Linear map and residual.}
% Keeping the linear part of \eqref{eq:dx-rec}--\eqref{eq:dg} and collecting the rest into
% a residual $R_{k}$, and using
% $\Delta x_{k+1} = (I-\alpha_{k}H)\Delta x_{k} + \alpha_{k}\beta_{k}\,\Delta d_{k-1}$,
% $\Delta d_{k} = -H\Delta x_{k} + \beta_{k}\Delta d_{k-1}$,
% \begin{equation}\label{eq:wrec}
% w_{k+1} = M_{k} w_{k} + R_{k},
% \qquad
% M_{k} = \begin{pmatrix} I - \alpha_{k} H & \alpha_{k}\beta_{k} I \\ -H & \beta_{k} I \end{pmatrix}.
% \end{equation}
% The map $M_{k}$ has the same eigenvalue structure as in the improve-or-localize
% analysis. The residual collects the three perturbing terms,
% \begin{equation}\label{eq:Rdef}
% R_{k} = (e_{k} - e_{k}') + \Delta\beta_{k}\, d_{k-1}' + \Delta\alpha_{k}\, d_{k}',
% \end{equation}
% whose first term satisfies, by \eqref{eq:ediff},
% $\norm{e_{k}-e_{k}'}\leq \Ltwo S\,\norm{\Delta x_{k}}\leq \Ltwo S\,\norm{w_{k}}$.
% % NOTE/TODO (open question from the notes): the stepsize-difference term Delta alpha_k.
% % Idea: set alpha_k = min{1/L1, C ||g_k||^{1+p-2q}} so the worst case (smallest step,
% % most likely to stay stuck) is selected, which should let us control Delta alpha_k.

% \paragraph{Bounding the $\beta$-difference (Fletcher--Reeves).}
% With $\beta^{\mathrm{FR}}_{k+1} = \norm{g_{k+1}}^{2}/\norm{g_{k}}^{2}$, splitting the
% numerator and denominator differences gives
% \[
% \Delta\beta_{k+1}
%   = \frac{\norm{g_{k+1}}^{2} - \norm{g_{k+1}'}^{2}}{\norm{g_{k}}^{2}}
%     \;-\; \beta_{k+1}'\,\frac{\norm{g_{k}}^{2} - \norm{g_{k}'}^{2}}{\norm{g_{k}}^{2}},
% \]
% with $\beta_{k+1}' = \norm{g_{k+1}'}^{2}/\norm{g_{k}'}^{2}$. For the numerators, with
% $\bar g$ an upper bound on the gradient norms ($\norm{g_{k+1}},\norm{g_{k+1}'}\leq\bar g$),
% the reverse triangle inequality and \Cref{ass:gradlip} give
% \[
% \big|\,\norm{g_{k+1}}^{2}-\norm{g_{k+1}'}^{2}\,\big|
%   = \big|\norm{g_{k+1}}-\norm{g_{k+1}'}\big|\big(\norm{g_{k+1}}+\norm{g_{k+1}'}\big)
%   \leq 2\bar g\,\norm{\Delta g_{k+1}} \leq 2\bar g\Lone\,\norm{\Delta x_{k+1}},
% \]
% and similarly $|\norm{g_{k}}^{2}-\norm{g_{k}'}^{2}|\leq 2\bar g\Lone\norm{\Delta x_{k}}$.
% For the denominator, let $\underline g$ satisfy $\norm{g_{k}}^{2}\geq\underline g^{2}$.
% With $|\beta_{k+1}'|\leq\bar\beta$ and $\norm{\Delta x_{k+1}}\leq\norm{w_{k+1}}$,
% $\norm{\Delta x_{k}}\leq\norm{w_{k}}$,
% \begin{equation}\label{eq:dbeta}
% |\Delta\beta_{k+1}|
%   \leq \frac{2\bar g\Lone}{\underline g^{2}}
%        \big(\norm{\Delta x_{k+1}} + |\beta_{k+1}'|\,\norm{\Delta x_{k}}\big)
%   \leq \frac{2\bar g\Lone}{\underline g^{2}}
%        \big(\norm{w_{k+1}} + \bar\beta\,\norm{w_{k}}\big).
% \end{equation}
% The corresponding residual term then obeys, using
% $\norm{d_{k}'}\leq\kappa\norm{g_{k}'}^{q}\leq\kappa\bar g^{\,q}$,
% \begin{equation}\label{eq:dbeta-d}
% \norm{\Delta\beta_{k+1}\, d_{k}'}
%   = |\Delta\beta_{k+1}|\,\norm{d_{k}'}
%   \leq \frac{2\bar g\Lone}{\underline g^{2}}
%        \big(\norm{w_{k+1}} + \bar\beta\,\norm{w_{k}}\big)\kappa\,\bar g^{\,q}.
% \end{equation}

% \paragraph{The gradient-norm constants.}
% While the iterate is stuck in the localization ball, $\norm{x_{k}-\xstar}\leq S$, so by
% \Cref{ass:gradlip} and $\grad f(\xstar)=0$,
% \[
% \norm{g_{k}} = \norm{\grad f(x_{k}) - \grad f(\xstar)} \leq \Lone\norm{x_{k}-\xstar}
%   \leq \Lone S,
% \qquad\text{hence}\qquad \bar g = \Lone S .
% \]
% For the lower bound $\underline g$: the perturbation is triggered because
% $\norm{g_{k}}\leq\epsg$, and at the post-perturbation point $g_{0} = \grad f(\snap+\xi)$\dots
% % TODO: finish the lower bound underline{g}; the notes stop here.

% \begin{remark}[Status of Step 3c]\label{rem:status3c}
% The coupled recursion \eqref{eq:wrec} and the residual bounds \eqref{eq:ediff} and
% \eqref{eq:dbeta-d} are in place. What remains: (i) finish the lower bound
% $\underline g$; (ii) settle the treatment of $\Delta\alpha_{k}$ (see the note after
% \eqref{eq:Rdef}); (iii) reconcile the $\beta$-index between $R_{k}$
% (term $\Delta\beta_{k} d_{k-1}'$) and the Fletcher--Reeves computation (done for
% $\Delta\beta_{k+1}$); and (iv) propagate $w_{k+1} = M_{k}w_{k} + R_{k}$ to show that
% $\norm{w_{k}}$ grows, then convert this into the high-probability escape and the final
% complexity bound.
% \end{remark}



\bibliographystyle{plain}
\bibliography{refs}

\end{document}